% MAS1004 Assignment 1: long report.
%
% You fill this in with your agent, in your repository. AGENTS.md tells the
% agent how. Red text in the PDF is what still has to be replaced.
%
% - Pictures go in figures/. Copy them from results/, which is not in git.
% - Outputs go in outputs/. run.py saves what it prints in
%   results/<tag>_output.txt, and PLAN.md saves the other outputs in results/
%   with `| tee results/...`. Copy them here. \outputfile{outputs/check.txt}
%   prints a file exactly as it is. Every number in this report has to come
%   from such an output.
%
% To make the PDF: compress this folder into a zip, and on Overleaf choose
% New Project, Upload Project. Set Menu, Compiler to XeLaTeX, and Recompile.

\documentclass[11pt]{article}
\usepackage{mas1004}

% Problem 7: three or four pictures in a row, each with a line of text under it.
\newcommand{\pic}[3][0.235]{%
  \begin{minipage}[t]{#1\linewidth}\centering
    \IfFileExists{#2}{\includegraphics[width=\linewidth,height=4.2cm,keepaspectratio]{#2}}{\picturefile{#2}}%
    \par\scriptsize #3
  \end{minipage}}

\reporttitle{Assignment 1: Long report}
\studentname{Lee Cheonhee}
\studentid{20251183}

\begin{document}
\makeheader

% ---------------------------------------------------------------------------
\section*{Problem 1. Categories and images}

\subsection*{Why these categories}
한국에 사는 여치과(Tettigoniidae) 곤충 중에서 흔히 발견되며, 내가 구분하지 못하는 다섯 종을 골랐다. 이는 일상 중 친구들이 나에게 "이건 어떤 풀벌레야?"라고 물어보면 제일 답하기 어려웠던 것이 여치였기 때문이다. 앞으로 친구들 앞에서 여치를 발견하면 여치과의 어떤 종인지 대답할 수 있도록 여치를 공부해보고 또 분류기의 성능도 시험해보고자 위 카테고리를 선택하였다. 
\begin{itemize}
\item \texttt{atlanticus\_brunneri} 좀날개여치
\item \texttt{gampsocleis\_sedakovii\_obscura} 여치
\item \texttt{gampsocleis\_ussuriensis} 긴날개여치
\item \texttt{paratlanticus\_ussuriensis} 갈색여치
\item \texttt{tettigonia\_ussuriana} 중베짱이
\end{itemize}
위 다섯 종은 한국에서 가장 흔히 발견되는 종이며, 개인적으로 길에서 발견하였을 때 여치과임을 알아보기 쉽지 않을 것이라 판단한 종들이다. 이미 내 눈에 여치과임이 명백한 종들은 친구들 앞에서 쉽게 대답할 수 있을 것이기에 주제 선정 취지와 맞지 않는다고 생각했다.

\subsection*{What I expected the model to use}
학습 전에는 모델이 여치의 색을 통해 크게 갈색 부류와 초록 부류를 나눌 것이라 생각하였으며, 이후 몸의 모양과 무늬(날개 길이, 날개가 몸을 덮는 모양, 몸통과 다리의 생김새, 아랫배의 색 등)의 미세한 차이로 다섯 종을 구별할 것이라고 생각했다. 종마다 내가 생각한 특징은 Problem 7의 각 종 분석 개요에 적었고, Problem 7에서 하나씩 검증했다.

\subsection*{Images per category}
첫 학습(\texttt{run})에 쓴 정리 전 사진은 304장이었다(\texttt{run\_output.txt}): 좀날개여치 59, 여치 50, 긴날개여치 61, 갈색여치 81, 중베짱이 53. 정리 후 최종 모델(\texttt{fixed\_v2})은 학습 227장(43, 39, 43, 59, 43), 고정 테스트 50장(종마다 10장)으로 학습·평가했고, 새 출처 사진은 종마다 5장, 모두 25장이다(\texttt{fixed\_v2\_output.txt}, \texttt{check.txt}). \texttt{check.py}는 여치의 학습 사진이 39장뿐이라는 경고를 냈다.

% ---------------------------------------------------------------------------
\section*{Problem 2. Preparing an image}

\sidebyside{figures/prepare_original.jpg}{다운로드한 사진 (gampsocleis\_ussuriensis/추가\_0020.jpg, 1280 $\times$ 720)}
           {figures/prepare_result.png}{prepare\_image가 만든 224 $\times$ 224}

\texttt{prepare\_image}는 짧은 변을 256으로 줄이고(256 $\times$ 455) 가운데 224 $\times$ 224만 남긴다. 이 사진은 그 과정에서 전체의 57\%가 버려졌다(\texttt{report\_extras.txt}). 가로로 긴 사진이라 양옆이 잘렸고, 오른쪽으로 뻗은 날개 끝과 다리가 잘려 나갔다. 긴날개여치는 날개 길이가 내가 생각한 가장 중요한 특징이었는데, 바로 그 날개 끝이 모델에게 보이지 않았다. 그래서 이 사진은 곤충이 가운데에 오도록 직접 자른 사본(\texttt{추가\_0020(cut).jpg})으로 바꿨다. 이렇게 직접 자른 사본으로 바꾼 사진은 모두 47장이다(Problem 4의 표).

% ---------------------------------------------------------------------------
\section*{Problem 3. First training runs}

\begin{center}\footnotesize
\begin{tabular}{llllrrrr}
\toprule
tag & start & trained & epochs & lr & trainable parameters & train acc. & test acc. \\
\midrule
run & ImageNet & all layers & 10 & 0.0001 & 11,179,077 & 100.0\% & 73.8\% \\
scratch & random & all layers & 10 & 0.0001 & 11,179,077 & 99.6\% & 31.1\% \\
frozen & ImageNet & last layer & 10 & 0.001 & 2,565 & 90.1\% & 63.9\% \\
custom & ImageNet & all layers & 4 & 0.0001 & 11,179,077 & 100.0\% & 78.7\% \\
\bottomrule
\end{tabular}
\end{center}
네 번 모두 정리 전 사진 304장을 무작위로 나눈 학습 243장, 테스트 61장으로 했다. 내 첫 번째 설정(\texttt{custom})은 에포크만 10에서 4로 줄인 것이다. 정리 뒤에는 테스트 사진을 \texttt{data/test}의 50장(종마다 10장)으로 고정하고 다시 학습했다. 이 표는 테스트 사진이 위 표와 달라서 서로 비교하지 않는다.

\begin{center}\footnotesize
\begin{tabular}{llllrrrr}
\toprule
tag & start & trained & epochs & lr & trainable parameters & train acc. & test acc. \\
\midrule
fixed & ImageNet & all layers & 10 & 0.0001 & 11,179,077 & 100.0\% & 88.0\% \\
scratch\_fixed & random & all layers & 10 & 0.0001 & 11,179,077 & 100.0\% & 30.0\% \\
frozen\_fixed & ImageNet & last layer & 10 & 0.0001 & 2,565 & 42.6\% & 36.0\% \\
frozen\_fixed\_lr3 & ImageNet & last layer & 20 & 0.001 & 2,565 & 95.2\% & 78.0\% \\
fixed\_v2 & ImageNet & all layers & 10 & 0.0001 & 11,179,077 & 100.0\% & 78.0\% \\
scratch\_fixed\_v2 & random & all layers & 10 & 0.0001 & 11,179,077 & 97.8\% & 32.0\% \\
frozen\_fixed\_v2 & ImageNet & last layer & 10 & 0.0001 & 2,565 & 41.0\% & 28.0\% \\
frozen\_fixed\_lr3\_v2 & ImageNet & last layer & 20 & 0.001 & 2,565 & 93.4\% & 78.0\% \\
\bottomrule
\end{tabular}
\end{center}
\texttt{fixed}는 학습 249장, \texttt{\_v2}는 정리를 더 한 뒤의 학습 227장이다. 웹 데모에 올린 최종 모델은 \texttt{fixed\_v2}다.

\texttt{frozen\_fixed\_lr3}는 내 두 번째 설정이다. 마지막 층만 학습한 \texttt{frozen\_fixed}(lr 0.0001, 10 에포크)는 학습 정확도가 0.221에서 0.422로, 테스트 정확도가 0.320에서 0.360으로 마지막 에포크까지 계속 올라가고 있었고 학습 손실도 1.6580에서 1.4032로 계속 내려가고 있었다(\texttt{frozen\_fixed\_output.txt}). 아직 더 올라갈 여지가 있다고 보고, 더 학습시키기 위해 학습률을 0.001로 올리고 에포크를 20으로 늘렸다. 그러자 테스트 정확도가 20 에포크 동안 0.280에서 0.780까지 올랐다(\texttt{frozen\_fixed\_lr3\_output.txt}). 정리 후의 사진(\texttt{\_v2})에도 같은 두 설정을 적용했다.

\picturefile{figures/run_curves.png}

손실은 내려갔다. 학습 손실은 1.2979(1 에포크)에서 4 에포크에 0.0546까지 빠르게 떨어진 뒤 0.01 근처에서 평평해졌다. 테스트 손실은 0.9039에서 6 에포크의 0.5525까지 내려간 뒤 더 내려가지 않고 0.5723으로 조금 올랐다. 테스트 정확도도 3--4 에포크의 0.803이 가장 높았고 마지막은 0.738이었다. 그래서 에포크를 4로 줄인 \texttt{custom}을 내 설정으로 해 보았고, 테스트 정확도가 78.7\%로 올랐다.

학습 정확도는 100.0\%, 테스트 정확도는 73.8\%로 26.2\%p 차이가 난다. 모델이 학습 사진은 모두 맞히지만 처음 보는 사진에서는 네 장 중 한 장 꼴로 틀린다는 뜻이다. 학습 사진의 종 특징뿐 아니라 사진마다의 배경이나 찍힌 모습까지 외웠다고 볼 수 있다(과적합). Problem 7에서 모델이 배경과 밝기를 단서로 쓴다는 것을 확인했는데, 이것도 같은 이야기다.

\begin{center}\picturefile[0.5\linewidth]{figures/run_confusion.png}\end{center}

가장 많이 섞인 두 종은 좀날개여치와 갈색여치다. 좀날개여치 4장이 갈색여치로, 갈색여치 3장이 좀날개여치로 분류되었다. 다음은 중베짱이 → 긴날개여치 3장, 여치 → 긴날개여치 3장이다. 좀날개여치와 갈색여치는 둘 다 갈색 계열이고 몸이 통통해서 어느 정도 예상했지만, 모양이 가장 닮았다고 생각한 긴날개여치와 중베짱이보다 더 많이 섞인 것은 의외였다. Problem 7에서 이 두 종을 배경(검정·풀잎)이 크게 갈라 놓는다는 것을 확인했다.

같은 사진, 같은 코드, 같은 에포크에서 ImageNet으로 시작하면 73.8\%, 무작위로 시작하면 31.1\%로 42.7\%p 차이가 났다(테스트 고정 후에도 88.0\% 대 30.0\%, 78.0\% 대 32.0\%). 무작위로 시작한 모델도 학습 정확도는 99.6\%였으므로 학습 사진은 외웠지만, 243장만으로는 일반적인 시각 특징(가장자리, 질감, 모양)부터 배울 수 없었다. ImageNet으로 시작한 모델은 120만 장에서 배운 특징을 그대로 쓰고 마지막 판단만 다섯 종에 맞게 바꾸면 되기 때문이다.

마지막 층만 학습하면 바뀌는 파라미터는 2,565개로 전체 11,179,077개의 0.02\%다. lr 0.001, 10 에포크에서 63.9\%로, 전부 학습한 73.8\%보다 9.9\%p 낮았다. 테스트를 고정한 뒤에는 lr 0.0001, 10 에포크에서 36.0\%(\_v2는 28.0\%)에 그쳤지만 정확도가 계속 오르는 중이었다. 그래서 두 번째 설정으로 lr 0.001에 20 에포크로 늘려 더 학습시키자 78.0\%로 전부 학습한 \texttt{fixed\_v2}(78.0\%)와 같아졌다. 마지막 층만으로도 충분히 따라잡지만, 학습률과 에포크를 더 줘야 했다.

\sidebyside{figures/fixed_v2_curves.png}{fixed\_v2 학습 곡선}
           {figures/fixed_v2_confusion.png}{fixed\_v2 혼동행렬 (고정 테스트 50장)}
\sidebyside{figures/scratch_curves.png}{scratch 학습 곡선}
           {figures/frozen_curves.png}{frozen 학습 곡선}

\subsection*{Output of the runs}
\outputfile{outputs/run_output.txt}
\outputfile{outputs/scratch_output.txt}
\outputfile{outputs/frozen_output.txt}
\outputfile{outputs/custom_output.txt}
\outputfile{outputs/fixed_output.txt}
\outputfile{outputs/scratch_fixed_output.txt}
\outputfile{outputs/frozen_fixed_output.txt}
\outputfile{outputs/frozen_fixed_lr3_output.txt}
\outputfile{outputs/fixed_v2_output.txt}
\outputfile{outputs/scratch_fixed_v2_output.txt}
\outputfile{outputs/frozen_fixed_v2_output.txt}
\outputfile{outputs/frozen_fixed_lr3_v2_output.txt}

% ---------------------------------------------------------------------------
\section*{Problem 4. Cleaning}

\subsection*{My rule}
분류기의 정확도를 올리기 위해서 각 종의 특징이 잘 드러난 이미지를 중심으로 학습시키고자 하였다. "배경색이 곤충과 너무 비슷하거나, 너무 어둡거나, 주위 나무나 풀이 많아서 곤충 위치 식별이 어려운 이미지"는 뺐다. 여기에 종을 확인하기 어려운 사진(다른 종, 약충, 메뚜기, 식별 특징이 안 보이는 구도, 낮은 해상도)도 직접 판별하여 뺐다. 
곤충이 작게 찍혔거나 가운데에 있지 않아 모델이 보는 가운데 정사각형 밖으로 잘리는 사진은 지우는 대신 곤충이 가운데 오도록 직접 자른 사본으로 바꿨다. 특히 긴날개여치의 경우 날개 끝이 화면에 다 들어오도록, 갈색여치의 경우 초록색인 아랫배가 화면에 들어오도록 잘랐다. 이는 분류기가 종의 주요 특징을 학습하게 만들기 위한 의도였다.
정리 전에 따로 적어 둔 규칙 문장은 다음과 같다. "곤충의 위치 식별이 어렵거나 쉐입을 파악하기 어려운 이미지는 제거한다. 잘못된 라벨의 이미지도 제거한다. 변환된 이미지에서 각 종의 특징이 잘릴 경우, 해당 특징이 화면 안에 들어오도록 조정한다."

\subsection*{Three kinds of junk}
\par\noindent
\pic[0.31]{figures/junk_1_grasshopper.jpg}{긴날개여치로 받은 메뚜기 (추가\_0014.jpg)}\hfill
\pic[0.31]{figures/junk_2_nymph.jpg}{종을 알 수 없는 약충 (2272-m0tkz.jpg)}\hfill
\pic[0.31]{figures/junk_3_too_dark.jpg}{너무 어두운 사진 (갈색여치 add\_025.jpg)}
\par\medskip
첫째, 다른 종이나 아예 여치가 아닌 곤충(메뚜기)이 그 종 이름으로 내려받아졌다. 둘째, 날개가 덜 자란 약충은 종을 알 수 없었다. 셋째, 너무 어둡거나 해상도가 낮아 특징이 보이지 않는 사진이 있었다.

\subsection*{What I removed}
\outputfile{outputs/clean_count.txt}

\texttt{downloaded}는 실제 다운로드 기록이 아니다. 다운로드 원본을 보존하지 못해서, \texttt{data/raw}를 \texttt{data/clean}과 \texttt{data/removed}로부터 다시 만들었다. \texttt{other}가 음수인 것은 직접 자른 사본 48장이 \texttt{data/clean}에는 있지만 \texttt{data/raw}에는 없기 때문이다. 제거 이유 중 가장 많은 47장은 사진을 지운 것이 아니라 직접 자른 사본으로 바꾼 것이다.

유사 복사본(유사도 0.95 초과)은 \texttt{clean.py suspects}가 한 쌍도 찾지 않았다. 의심 목록을 만든 249장으로 다시 재어도 0쌍이었다(\texttt{report\_extras.txt}). 대신 여치 한 장(\texttt{다운로드 (5).jpg})이 고정 테스트의 사진과 똑같은 파일이어서 학습 사진에서 뺐다.

종마다 의심 목록 상위 20장 중 제거한 것은 좀날개여치 7장, 여치 1장, 긴날개여치 9장, 갈색여치 6장, 중베짱이 4장, 모두 27장이다. 그중 13장은 다른 종·약충·메뚜기·어두움·낮은 해상도·구도처럼 규칙에 어긋나 지운 것이고, 14장은 직접 자른 사본으로 바꾼 것이다(\texttt{report\_extras.txt}). 제거한 81장 중 54장은 상위 20장에 없었으므로, 의심 목록에 없는 잘못된 사진도 많았다. 특히 종 이름이 틀린 사진 10장은 테스트 사진을 검토하다가 찾았다.

나는 직접 사진을 보면서 정제/제거 대상 사진 13장을 먼저 골랐고, 테스트 사진을 검토하면서 11장을 더 골랐다. 그 뒤에 의심 목록을 보았다. 눈으로 고른 사진 중 다수는 의심 목록 상위에도 있었다(예: 좀날개여치 추가\_0018은 2위). 의심 목록을 본 뒤 정크 이미지라는 걸 알아차린 경우도 있었다. 상위 20장에서 중베짱이 add\_045(식별 특징이 안 보이는 구도), 갈색여치 add\_034(다른 종), add\_024(낮은 해상도) 3장은 직접 볼 때 놓쳤다가 의심 목록을 보고 알아차려 제거하였다.

\subsection*{Before and after cleaning}
정리는 두 단계로 했다(\texttt{data/removed/log.csv}).
\begin{itemize}
\item 1차 정리: 고정 테스트 50장(종마다 10장)을 \texttt{data/test}로 떼어 내면서 테스트 사진을 검토했고, 그때 찾은 종 라벨이 틀린 사진 10장과 테스트 사진과 똑같은 중복 1장, 모두 11장을 뺐다(좀날개여치 3, 여치 4, 긴날개여치 1, 갈색여치 3). 학습 249장으로 \texttt{fixed}를 학습했다.
\item 2차 정리: 규칙에 따라 23장을 더 뺐고(좀날개여치 6, 긴날개여치 7, 갈색여치 9, 중베짱이 1), 곤충이 가운데 오지 않아 잘리던 47장은 직접 자른 사본으로 바꿨다(좀날개여치 8, 여치 4, 긴날개여치 16, 갈색여치 5, 중베짱이 14). 학습 227장으로 \texttt{fixed\_v2}를 학습했다.
\end{itemize}

\outputfile{outputs/compare.txt}

\begin{center}\footnotesize
\begin{tabular}{lrr}
\toprule
모델 & 테스트 정확도 & 새 출처 사진 \\
\midrule
run (정리 전) & 73.8\% (무작위 61장) & 72.0\% (18/25) \\
fixed (1차 정리) & 88.0\% (고정 50장) & 76.0\% (19/25) \\
fixed\_v2 (2차 정리, 최종) & 78.0\% (고정 50장) & 72.0\% (18/25) \\
\bottomrule
\end{tabular}
\end{center}

1차 정리의 경우, 정리하면서 테스트 사진도 바뀌었으므로(무작위 61장 → 고정 50장), 테스트 정확도는 정리 전후를 직접 비교할 수 없다. 바뀌지 않는 새 출처 사진 25장에서는 정리 전 72.0\%, 1차 정리 후 76.0\%, 2차 정리 후 72.0\%로, 정리가 새 사진의 정확도를 거의 바꾸지 못했다. 같은 고정 테스트에서도 2차 정리 뒤 88.0\%에서 78.0\%로 떨어졌다. 학습 설정은 같고 2차 정리(23장 제거, 47장을 직접 자른 사본으로 교체)만 달라서, 원인은 확인하지 못했다. 다만 Problem 7에서 모델이 곤충의 모양보다 배경, 곤충의 밝기, 날개가 몸을 덮은 인상에 기댄다는 것을 확인했다. 라벨이 틀린 사진을 빼고 곤충이 가운데 오도록 자르는 정리로는 이런 단서가 바뀌지 않기 때문에 새 사진의 정확도가 그대로였다고 생각한다. 테스트 50장에서 한 장은 2\%p, 새 사진 25장에서 한 장은 4\%p라서 작은 차이는 우연일 수도 있다.

% ---------------------------------------------------------------------------
\section*{Problem 5. Images from a new source}

새 출처 사진 25장(종마다 5장, \texttt{data/new\_images}의 p1--p5)은 대부분 `곤충나라 식물나라'라는 비공개 네이버 카페에서 구했다. 이 카페는 회원이 직접 발견하거나 키우는 곤충의 사진을 올리는 곳이고 비공개라서, 학습 데이터와 겹칠 일이 없다. 나머지 소수의 사진은 네이버 블로그의 이미지다. 학습 사진은 구글 이미지 검색으로 내려받았기 때문에 블로그 이미지와 겹칠 가능성은 적지만, 블로그 이미지는 검색에 나올 수 있어 겹치지 않는다고 보장할 수는 없다. 그래서 아래의 유사도 검사로 새 사진이 학습 데이터와 다르다는 것을 확인했다.

\outputfile{outputs/overlap.txt}
새 사진 25장 중 다운로드한 308장과 유사 복사본인 것은 없었다.

\begin{center}\footnotesize
\begin{tabular}{lrr}
\toprule
& 고정 테스트 (다운로드) & 새 출처 사진 \\
\midrule
fixed\_v2 정확도 & 78.0\% (50장) & 72.0\% (18/25) \\
\bottomrule
\end{tabular}
\end{center}

\picturefile{figures/wrong_new_images.png}

\texttt{check.py}가 꼽은 가장 확신한 오답 5장과, 무엇이 모델을 틀리게 밀었다고 생각하는지다. 원인은 Problem 7에서 이미지를 바꿔 확인한 결과를 근거로 적었다.
\begin{enumerate}
\item 중베짱이 p5 → 긴날개여치 77\%: 위에서 찍은 사진이다. 모델이 중베짱이를 알아보는 단서인 `날개가 몸을 감싼 옆모습'이 보이지 않고, 위에서 본 모습은 긴날개여치 쪽으로 학습되어 있었다(가설 11, 13).
\item 긴날개여치 p4 → 중베짱이 75\%: 자갈 위에서 찍혔다. 배경만 회색으로 바꾸자 5.4\%에서 83.8\%로 정답이 되었다(가설 16). 모델은 회색 돌·콘크리트 배경을 긴날개여치의 단서로 배웠다.
\item 긴날개여치 p5 → 중베짱이 69\%: 곤충이 밝게 찍혔다. 곤충만 어둡게 하자 22.9\%에서 44.1\%로 정답이 되었다(가설 16).
\item 중베짱이 p4 → 좀날개여치 67\%: 위에서 찍어 등의 주황·갈색 줄무늬가 도드라졌다. 그 줄무늬만 초록으로 칠하자 14.7\%에서 64.1\%로 정답이 되었다(가설 11).
\item 긴날개여치 p2 → 중베짱이 66\%: 날개 아래로 드러난 배가 중베짱이 쪽 단서로 읽히고, 곤충도 밝다. 날개와 배의 경계를 가리면 73.6\%, 배만 가리면 65.0\%, 곤충만 어둡게 하면 65.3\%로 정답이 되었다(가설 14, 15, 16).
\end{enumerate}

새 사진이 다운로드 사진과 다른 점은 다음과 같다. 위에서 내려다보고 찍은 사진이 많다(중베짱이 p3--p5, 긴날개여치 p3). 배경이 자갈, 흰 판, 방충망, 손, 실내 바닥처럼 다운로드 사진에 흔한 회색 돌·콘크리트나 풀잎과 다르다. 조명도 강한 빛이나 플래시로 곤충이 밝게 찍히거나(긴날개여치 p2, p5), 반대로 곤충만 어둡게 찍힌 경우(중베짱이 p1)가 있다. 곤충의 크기와 거리도 사진마다 다르다. Problem 7에서 모델이 바로 이런 배경과 밝기를 단서로 쓴다는 것을 확인했기 때문에, 새 사진에서 정확도가 떨어진 것은 이 차이 때문이라고 생각한다.

% ---------------------------------------------------------------------------
\section*{Problem 6. The demo}

Address: \url{https://chunhee674.github.io/mas1004-assignment1/}

다섯 종이 아닌 것을 보여 줘도 모델은 반드시 다섯 종 중 하나로 답한다. `모르겠다'는 답이 없기 때문이다. 긴날개여치로 잘못 내려받았던 메뚜기 두 장을 넣자 둘 다 중베짱이로 답했다(46.7\%, 51.1\%, \texttt{report\_extras.txt}). 확률이 절반 정도로 높지 않았지만, \texttt{check.py}가 꼽은 확신한 오답들(62--77\%)과 비슷한 수준이라 확률만 보고 범주 밖의 사진이라고 알 수는 없다.

\sidebyside{figures/none_1.jpg}{메뚜기 추가\_0014: 중베짱이 46.7\%}
           {figures/none_2.jpg}{메뚜기 추가\_0022: 중베짱이 51.1\%}

\outputfile{outputs/report_extras.txt}

% ---------------------------------------------------------------------------
\section*{Problem 7. What the model looks at}

모든 수치는 정답 클래스의 확률이다. 좀날개여치 이미지는 좀날개여치일 확률, 갈색여치 이미지는 갈색여치일 확률을 적었다. 확률이 가장 높은 클래스가 정답이 아니면 오답으로 본다. 검증에는 학습에 쓰지 않은 이미지(test, new\_images)만 사용했다.

이미지 정제는 \texttt{results/problem7/}의 도구로 했고(2026-10-16), 정제한 이미지는 \texttt{data/changed/<종>/<처리>/}에 있다. 확률은 \texttt{results/problem7/measure.py}가 정제한 이미지 파일을 데모 페이지와 같은 방식(\texttt{docs/model.onnx}, \texttt{src/check.py}의 \texttt{prepare})으로 분류기에 넣어 출력한 값이고, 출력 전체는 이 절 끝에 있다. 배경 제거는 rembg의 birefnet-general 모델(\texttt{make\_masks.py})로 했다. 각 가설에서는 내가 세운 가설, 그것을 검증하려고 요청한 정제(``정제 요청'', 요청 없이 만든 것은 ``정제'')와 그때 에이전트에게 한 말, 그 결과와 해석을 차례로 적었다. 그림 아래의 ``정제N''은 본문에서 그 그림을 가리키는 번호다.

% ----- 좀날개여치
\subsection*{좀날개여치 오분류 원인 분석}

여치 5종 분류기(ResNet18, fixed\_v2)의 테스트 이미지 중 좀날개여치 10장 가운데 단 1장(test/add\_031)만 갈색여치로 오분류되었다. 이 이미지 하나를 대상으로 오답의 원인에 대한 가설을 세우고, 이미지를 정제해 다시 분류기에 넣는 방식으로 가설을 검증했다.

\subsection*{가설 1. 붉은 색감 때문에 갈색여치로 보았다}

오답 이미지(정제1)는 다른 좀날개여치 이미지들에 비해 여치의 몸이 붉은기가 돌도록 찍혀 있었다. 좀날개여치는 대부분 마른 잎 같은 밝은 갈색이므로, 붉은 색감이 갈색여치로 판단하게 만든 원인이라고 가정했다.

이를 확인하기 위해 붉은기를 없애고 몸색을 더 갈색에 가깝게 보정한 이미지(정제2)를 만들어 다시 분류기에 넣었다. 보정은 내가 직접 제미나이를 활용하여 했다(\texttt{data/changed/atlanticus\_brunneri/h1\_redfix\_by\_user/}).

에이전트에게 한 말 그대로:
\begin{quote}
""
\end{quote}

\sidebyside{figures/p7_001.jpg}{정제1 원본: 12.9\% (오답, 갈색여치)}
           {figures/p7_002.jpg}{정제2 붉은기 보정: 4.6\% (오답, 갈색여치)}

좀날개여치 정답 확률은 12.9\%에서 4.6\%로 오히려 낮아졌고, 여전히 갈색여치로 분류되었다. 따라서 붉은 색감은 오답의 원인이 아니었다.

\subsection*{가설 2. 풀잎 배경 때문에 갈색여치로 보았다}

내 눈에는 붉은기 외에 이 여치를 갈색여치로 볼 만한 근거가 도저히 보이지 않았다. 다리에 가시가 많지도 않고, 다리 색도 갈색여치와 달랐으며, 가슴의 문양과 몸 전체의 얼룩도 갈색여치와 달랐다. 그래서 여치가 아닌 부분, 즉 풀잎 배경이 원인이라고 가정했다.

배경 제거 도구(rembg의 birefnet-general 모델)로 원본(정제1)과 보정본(정제2)의 배경을 지우고 투명 PNG로 분류기에 넣었다.

\textbf{정제.} 배경을 지운 이미지는 처음에 내가 직접 만들어 넣어 보았다. 그런데 다리와 배까지 너무 많이 지워져서, 더 섬세한 배경 제거를 요청했다(birefnet-general). 같은 두 이미지의 지운 자리를 흰색으로 채운 이미지도 만들었다.

에이전트에게 한 말 그대로:
\begin{quote}
"아무리 생각해도 이상해. 내 눈으로 보기에는 좀날개여치가 명확하거든 혹시 배경 탓이 아닐까? 배경을 보고 좀날개여치와 갈색여치를 구분하는 걸수도 있어. 배경을 지워봐."
\end{quote}
\begin{quote}
"좋아 근데 아쉬운점이 있어. 배경을 제거할때 다리나 배가 너무 많이 날아가. 좀 더 섬세히 배경 제거는 못하니?"
\end{quote}

좀날개여치 정답 확률은 원본이 12.9\%에서 99.2\%(정제3), 보정본이 4.6\%에서 97.9\%(정제4)로 올라 모두 정답이 되었다.

그런데 투명 PNG는 분류기에 들어갈 때 투명한 부분이 검정으로 바뀐다. 즉 정제3·4는 배경을 없앤 것이 아니라 배경을 검정으로 바꾼 것과 같다. 정답이 된 이유가 풀잎이 사라져서인지, 검정 배경이 생겨서인지 구분하기 위해 같은 이미지의 지운 자리를 흰색으로 채워 다시 넣었다.

\par\noindent
\pic{figures/p7_003.jpg}{정제3 원본, 배경 제거(투명): 99.2\%}\hfill
\pic{figures/p7_004.jpg}{정제4 보정본, 배경 제거(투명): 97.9\%}\hfill
\pic{figures/p7_005.jpg}{정제5 원본, 흰 배경: 52.5\%}\hfill
\pic{figures/p7_006.jpg}{정제6 보정본, 흰 배경: 46.8\% (오답, 갈색여치)}
\par\medskip

흰 배경에서는 원본이 52.5\%(정제5)로 겨우 정답이 되었고, 보정본은 46.8\%(정제6)로 여전히 갈색여치로 분류되었다. 풀잎을 지우는 것만으로 정답 확률이 약 40\%p 올랐으므로, 이 이미지에서는 풀잎 배경이 오답에 영향을 준 것이 맞다.

하지만 흰 배경에서는 정답과 오답의 경계(50\% 안팎)에 머물렀고, 99\%의 확신을 만든 것은 검정 배경이었다. 따라서 `검정 배경이 분류를 좀날개여치 쪽으로 끌어당긴다'는 새로운 가설을 세웠다.

\subsection*{가설 3. 검정 배경이 분류를 좀날개여치 쪽으로 끌어당긴다}

이 가설이 맞다면 두 가지가 함께 나타나야 한다. 첫째, 풀잎 배경을 흰색으로 바꿔도 다른 이미지들의 정답은 크게 달라지지 않아야 한다. 둘째, 검정으로 바꾸면 좀날개여치 이미지는 정답 확률이 오르고 갈색여치 이미지는 정답 확률이 떨어져야 한다.

학습에 쓰지 않은 이미지 중 풀잎 배경인 좀날개여치 2장과 갈색여치 4장을 모아, 각각 배경을 제거한 뒤 흰색과 검정으로 채워 분류기에 넣었다.

\textbf{정제.} test와 new\_images에서 풀잎 배경인 좀날개여치 2장(test/add\_011, new\_images/p3)과 갈색여치 4장(test/다운로드 (1), test/추가\_0038, test/추가\_0050, new\_images/p3)의 배경을 지우고, 각각 흰색과 검정으로 채운 이미지를 만들었다.

에이전트에게 한 말 그대로:
\begin{quote}
"배경이 풀잎인 좀날개여치와 갈색여치 이미지를 찾고, 그 배경을 제거하고, 배경 전체를 흰색과 검정색 두 경우로 채워줘."
\end{quote}

\begin{center}\footnotesize
\begin{tabular}{lrrr}
\toprule
이미지 & 원본 & 흰 배경 & 검정 배경 \\
\midrule
좀날개여치 A (test/add\_011, 정제7--9) & 90.3\% & 81.8\% & 98.3\% \\
좀날개여치 B (new\_images/p3) & 92.3\% & 94.7\% & 98.2\% \\
갈색여치 A (test/다운로드 (1), 정제10--12) & 85.9\% & 44.2\% & 0.7\% (오답) \\
갈색여치 B (test/추가\_0038, 정제13--15) & 99.5\% & 74.0\% & 15.7\% (오답) \\
갈색여치 C (test/추가\_0050) & 98.1\% & 94.0\% & 37.0\% \\
갈색여치 D (new\_images/p3) & 98.6\% & 92.5\% & 24.4\% (오답) \\
\bottomrule
\end{tabular}
\end{center}

흰 배경에서는 6장 모두 정답을 유지했다. 다른 이미지에서는 풀잎을 지워도 분류가 바뀌지 않았으므로, 분류기가 `풀잎 배경이면 갈색여치'라는 일반적인 규칙을 배운 것은 아니었다.

검정 배경에서는 좀날개여치 2장의 정답 확률이 모두 올랐고(90.3\%→98.3\%, 92.3\%→98.2\%), 갈색여치 4장은 모두 크게 떨어져 3장이 좀날개여치로 오분류되었다. 두 예측이 모두 맞아 가설 3을 지지한다.

\par\noindent
\pic[0.31]{figures/p7_007.jpg}{정제7 좀날개여치 A 원본: 90.3\%}\hfill
\pic[0.31]{figures/p7_008.jpg}{정제8 흰 배경: 81.8\%}\hfill
\pic[0.31]{figures/p7_009.jpg}{정제9 검정 배경: 98.3\%}
\par\medskip
\par\noindent
\pic[0.31]{figures/p7_010.jpg}{정제10 갈색여치 A 원본: 85.9\%}\hfill
\pic[0.31]{figures/p7_011.jpg}{정제11 흰 배경: 44.2\%}\hfill
\pic[0.31]{figures/p7_012.jpg}{정제12 검정 배경: 0.7\% (오답, 좀날개여치)}
\par\medskip
\par\noindent
\pic[0.31]{figures/p7_013.jpg}{정제13 갈색여치 B 원본: 99.5\%}\hfill
\pic[0.31]{figures/p7_014.jpg}{정제14 흰 배경: 74.0\%}\hfill
\pic[0.31]{figures/p7_015.jpg}{정제15 검정 배경: 15.7\% (오답, 좀날개여치)}
\par\medskip

왜 검정 배경이 좀날개여치 쪽으로 기울었는지 학습 데이터를 확인했다(\texttt{results/problem7/train\_backgrounds.py}). 배경이 어두운 학습 이미지(가장자리 평균 밝기 70/255 미만)는 좀날개여치가 43장 중 5장(참고1--5)이었고, 갈색여치는 59장 중 0장이었다. 학습 데이터에서 어두운 배경은 좀날개여치에만 있었기 때문에, 분류기가 어두운 배경을 좀날개여치의 단서로 배운 것으로 보인다.

\par\noindent
\pic[0.19]{figures/p7_ref_1.jpg}{참고1}\hfill
\pic[0.19]{figures/p7_ref_2.jpg}{참고2}\hfill
\pic[0.19]{figures/p7_ref_3.jpg}{참고3}\hfill
\pic[0.19]{figures/p7_ref_4.jpg}{참고4}\hfill
\pic[0.19]{figures/p7_ref_5.jpg}{참고5}
\par\medskip

% ----- 갈색여치
\subsection*{갈색여치 분석 개요}

테스트와 new\_images의 갈색여치 15장 중 13장은 정답, 2장(test/add\_031, test/natureobserve)은 좀날개여치로 오분류되었다. 정답이었던 것 중에도 2장(test/add\_008 50.8\%, test/추가\_0039 46.7\%)은 간신히 1위를 한 경계 사례였다.

나는 갈색여치를 다음 특징으로 구별한다고 생각했다.
\begin{enumerate}
\item 초록색 아랫배
\item 검은 몸통
\item 흰색 혹은 주황색으로 얼룩덜룩한 가슴
\item 다리의 노란 중간 마디
\item 다리에 난 검은 가시
\end{enumerate}

오답 이미지 두 장을 보고, 나는 각각 다음을 오답의 원인으로 생각했다.
\begin{itemize}
\item test/add\_031: 자연광 때문에 가슴과 몸통이 평상시보다 밝게 찍혔고, 아랫배의 초록색이 거의 보이지 않는다.
\item test/natureobserve: 몸이 다른 갈색여치 이미지보다 밝고, 다리의 무늬와 가시가 뚜렷하게 나타나지 않았으며, 아랫배의 초록색이 잘 보이지 않는다.
\end{itemize}

이 짐작을 검증하려면 먼저 모델이 내 특징들을 실제로 쓰는지부터 알아야 한다. 이에 다음 순서로 검증했다.
\begin{itemize}
\item 가설 4: 15장 전체에서 특징을 하나씩 없애, 모델이 어느 특징에 기대는지 확인했다.
\item 가설 5: 그 결과를 바탕으로 오답 두 장에서 내가 짐작한 원인을 하나씩 검증했다.
\item 가설 6: 좀날개여치 분석에서 확인한 배경의 영향을 요소로 추가해, 몸통·가슴·다리와 교차해 비교했다.
\end{itemize}
수치는 모두 갈색여치 정답 확률이다.

\subsection*{가설 4. 모델은 내가 세운 특징으로 갈색여치를 구별한다}

어떤 특징을 없앴을 때 정답 확률이 크게 떨어지면 모델이 그 특징을 쓰고 있는 것이고, 그대로라면 쓰지 않는 것이다. 초록 아랫배와 검은 몸통은 색에 관한 특징이라 그 색을 바꿨다. 가슴과 다리는 무늬와 모양에 관한 특징이다. 색만 바꾸면 처리가 깔끔하지 않고 갈색여치 중에는 원래 다리가 검은 개체도 있어서, 부위를 통째로 가리거나 지웠다.
\begin{itemize}
\item 초록 아랫배 제거(특징 1): 곤충 안의 초록색을 갈색으로 바꿨다.
\item 검은 몸통 → 갈색(특징 2): 몸통(가슴과 배)의 검은 부분을 밝은 갈색으로 바꿨다.
\item 가슴 가림(특징 3): 머리 바로 뒤의 가슴(앞가슴등판)만 회색 네모로 가렸다. 네모는 사진마다 가슴 위치에 맞게 직접 정했다.
\item 다리 제거(특징 4·5): 다리를 지우고 그 자리를 주변 배경으로 메웠다. 노란 마디와 검은 가시가 함께 사라진다.
\end{itemize}
곤충 영역은 배경 제거 도구(birefnet)가 찾은 영역을 썼고, 그 안의 두꺼운 부분을 몸통, 가는 부분을 다리로 나눴다. 갈색여치 15장 모두에 네 처리를 각각 따로 했다(\texttt{change\_paratlanticus.py}).

\textbf{정제 요청.} 갈색여치 15장 각각에 대해 네 가지 이미지를 요청했다. (1) 곤충 안의 초록색을 갈색으로 바꾼 이미지, (2) 몸통의 검은 부분을 밝은 갈색으로 바꾼 이미지, (3) 머리 바로 뒤의 가슴(앞가슴등판)만 회색 네모로 가린 이미지, (4) 다리를 지우고 그 자리를 주변 배경으로 메운 이미지다. 처음에는 다리의 검은 가시를 지우고 노란 마디를 칠하는 방법을 썼지만, 가시는 깔끔하게 지워지지 않고 마디를 칠하면 다리가 검어지는데 갈색여치 중에는 원래 다리가 검은 개체도 있어서, 다리를 통째로 지우는 것으로 바꿨다. 가슴 네모도 처음에는 몸통 전체를 덮어서, 머리 아래 가슴만 덮도록 다시 정하게 했다.

에이전트에게 한 말 그대로:
\begin{quote}
"정제할 이미지 종류는 총 4개야. 1. 갈색여치의 아랫배를 보면 초록색이야. 이걸 갈색으로 바꿔줘. 2. 몸통을 좀 더 밝은 갈색으로 보정해줘. 3.머리 바로 뒤의 가슴을 보면 무늬가 있어. 이 부분을 가릴거야. 머리 바로 뒤 가슴의 옆부분에도 무늬가 있으니까 갈색여치의 옆모습을 찍은 사진이든 위에서 찍은 사진이든 가슴의 무늬가 보이지 않도록 박스로 가려야해. 4.다리의 중간 마디에 노란색 무늬가 있어. 이걸 칠해서 가려줘. 5.다리를 보면 검정 가시가 있어. 가시를 지워줘."
\end{quote}
\begin{quote}
"다리의 검은 가시 없애는 것과 노란 마디를 칠하는 건 너무 애매한 방식이야. 검은 가시는 깔끔하게 없어지지 않고 있고 노란 마디를 칠하면 검은 다리가 되는데 이게 갈색 여치 중에 검은 다리를 가진 애들도 있거든. 그러니 그냥 다리를 제거해서 검증하자. 다리를 제거했는데 영향이 있으면 다리의 가시와 무늬 중에 영향을 미치는 요소가 있다고 하면 되잖아."
\end{quote}

\begin{center}\footnotesize
\begin{tabular}{lrrr}
\toprule
없앤 요소 & 정답 확률 평균 변화 (정답이던 13장) & 오답으로 바뀜 & 15장 중 정답 \\
\midrule
초록 아랫배 (특징 1) & +0.5\%p & 0/13 & 13/15 \\
검은 몸통 (특징 2) & $-$33.0\%p & 6/13 & 7/15 \\
가슴 (특징 3) & $-$11.6\%p & 2/13 & 12/15 \\
다리 (특징 4·5) & $-$35.4\%p & 5/13 & 8/15 \\
\bottomrule
\end{tabular}
\end{center}

초록 아랫배는 쓰이지 않았다. 아랫배의 초록이 뚜렷한 new\_images/p1도 초록을 갈색으로 바꾸기 전후가 86.8\%(정제16)와 87.3\%(정제17)로 거의 같았다.

\sidebyside{figures/p7_016.jpg}{정제16 new\_images/p1 원본: 86.8\%}
           {figures/p7_017.jpg}{정제17 초록 아랫배 제거: 87.3\%}

검은 몸통은 큰 단서였다. 몸통을 갈색으로 바꾸자 6장이 오답이 되었다(5장은 좀날개여치, 1장은 긴날개여치). new\_images/p4는 72.3\%(정제18)에서 4.5\%(정제19)로, test/add\_013은 90.2\%(정제20)에서 4.8\%(정제21)로 떨어졌다.

\sidebyside{figures/p7_018.jpg}{정제18 new\_images/p4 원본: 72.3\%}
           {figures/p7_019.jpg}{정제19 검은 몸통 → 갈색: 4.5\% (오답, 좀날개여치)}
\sidebyside{figures/p7_020.jpg}{정제20 test/add\_013 원본: 90.2\%}
           {figures/p7_021.jpg}{정제21 검은 몸통 → 갈색: 4.8\% (오답, 좀날개여치)}

가슴 무늬는 작은 단서였다. 가슴을 가리면 정답 확률이 낮아지기는 했지만(p4: 72.3\% → 47.5\%, 정제22) 오답이 된 것은 2장뿐이었다. 오히려 오답이던 test/add\_031은 가슴을 가리자 34.7\%(정제23)에서 54.4\%(정제24)로 올라 정답이 되었다. 이 개체의 밝은 주황빛 가슴이 분류를 좀날개여치 쪽으로 끌고 있었다는 뜻이다.

\sidebyside{figures/p7_018.jpg}{정제18 new\_images/p4 원본: 72.3\%}
           {figures/p7_022.jpg}{정제22 가슴 가림: 47.5\%}
\sidebyside{figures/p7_023.jpg}{정제23 test/add\_031 원본: 34.7\% (오답, 좀날개여치)}
           {figures/p7_024.jpg}{정제24 가슴 가림: 54.4\%}

다리는 가장 큰 단서였다. 다리를 지우자 5장이 오답이 되었다. p4는 72.3\%(정제18)에서 6.4\%(정제25)로, test/add\_023은 90.4\%(정제26)에서 8.7\%(정제27)로 떨어졌다. 다리를 통째로 지웠으므로 가시와 무늬 중 무엇이 영향을 주었는지까지는 나눌 수 없지만, 그중 모델이 쓰는 요소가 있다는 것은 확인했다.

\sidebyside{figures/p7_018.jpg}{정제18 new\_images/p4 원본: 72.3\%}
           {figures/p7_025.jpg}{정제25 다리 제거: 6.4\% (오답)}
\sidebyside{figures/p7_026.jpg}{정제26 test/add\_023 원본: 90.4\%}
           {figures/p7_027.jpg}{정제27 다리 제거: 8.7\% (오답)}

정리하면 내 특징 중 모델이 실제로 기대는 것은 검은 몸통과 다리이고, 가슴 무늬는 약하게, 초록 아랫배는 전혀 쓰지 않았다. 이미지별 전체 결과는 다음과 같다.

\begin{center}\footnotesize
\begin{tabular}{lrrrrr}
\toprule
이미지 & 원본 & 초록 아랫배 & 검은 몸통 & 가슴 가림 & 다리 제거 \\
\midrule
test/add\_008 & 50.8\% & 50.8\% & 14.1\% (오답) & 63.2\% & 16.4\% (오답) \\
test/add\_013 & 90.2\% & 90.7\% & 4.8\% (오답) & 55.1\% & 54.0\% \\
test/add\_023 & 90.4\% & 90.4\% & 34.6\% (오답) & 93.3\% & 8.7\% (오답) \\
test/add\_031 & 34.7\% (오답) & 34.8\% (오답) & 1.1\% (오답) & 54.4\% & 13.8\% (오답) \\
test/natureobserve & 18.9\% (오답) & 18.6\% (오답) & 0.7\% (오답) & 4.9\% (오답) & 8.2\% (오답) \\
test/다운로드 (1) & 85.9\% & 83.7\% & 78.8\% & 63.2\% & 28.4\% (오답) \\
test/추가\_0005 & 68.2\% & 71.5\% & 15.0\% (오답) & 31.5\% (오답) & 10.4\% (오답) \\
test/추가\_0038 & 99.5\% & 98.6\% & 70.7\% & 99.0\% & 92.2\% \\
test/추가\_0039 & 46.7\% & 49.4\% & 7.3\% (오답) & 28.0\% (오답) & 34.9\% \\
test/추가\_0050 & 98.1\% & 97.9\% & 90.6\% & 97.1\% & 96.4\% \\
new/p1 & 86.8\% & 87.3\% & 83.1\% & 75.1\% & 53.6\% \\
new/p2 & 94.5\% & 96.6\% & 75.1\% & 93.9\% & 41.1\% \\
new/p3 & 98.6\% & 99.1\% & 76.2\% & 94.0\% & 93.9\% \\
new/p4 & 72.3\% & 71.8\% & 4.5\% (오답) & 47.5\% & 6.4\% (오답) \\
new/p5 & 95.1\% & 96.0\% & 93.7\% & 85.2\% & 80.1\% \\
\bottomrule
\end{tabular}
\end{center}

\subsection*{가설 5. 오답 두 장은 내가 짐작한 원인 때문에 틀렸다}

가설 4에서 모델이 검은 몸통과 다리를 크게 쓰고 아랫배 초록은 쓰지 않는다는 것을 확인했다. 이를 바탕으로 오답 두 장에 대해 처음 짐작한 세 가지 원인을 하나씩 검증했다.

아랫배의 초록이 안 보임 (두 장). 가설 4에서 모델은 아랫배 초록을 전혀 쓰지 않았다. 따라서 초록이 안 보인 것은 오답의 원인이 아니다.

몸통이 밝게 찍힘 (두 장). 몸통이 밝은 것이 원인이라면, 몸통을 어둡게 했을 때 정답이 되어야 한다. 처음에는 몸통 전체의 밝기를 낮췄는데, 흰색·주황 무늬까지 함께 어두워져 무늬가 사라졌다. 실제 갈색여치는 검은 바탕 위에 밝은 무늬가 선명하므로, 몸통의 어두운 바탕 부분만 검게 하고 밝은 무늬는 원래 색대로 남겨 다시 만들었다.

\textbf{정제 요청.} 몸통의 어두운 바탕만 검게 하고 흰색·주황 무늬는 원래 색대로 남긴 이미지를 오답 두 장과 갈색여치 15장 전체에 대해 요청했다. 처음 만든 이미지는 몸통 전체의 밝기를 낮춘 것이라 무늬까지 사라졌기 때문이다. natureobserve의 흐린 다리는 다리를 지운 이미지(가설 4)로 검증했다.

에이전트에게 한 말 그대로:
\begin{quote}
"갈색 여치의 오답 두장에서 몸통 밝기를 낮추는 정제와 다리를 지우는 정제해줘."
\end{quote}
\begin{quote}
"몸통을 검게 만드는 거에서 내 눈에 보기엔 검은 칠 한게 몸의 무늬를 다 상쇄시키는 것처럼 보여. 몸의 무늬 정보가 많이 손실되는 것 같아. 몸의 무늬 정보는 남도록 무늬는 건들지 말고 무늬 이외의 몸통 색 밝기를 낮춰줘."
\end{quote}

그러자 add\_031은 34.7\%(정제23)에서 67.9\%(정제28)로, natureobserve는 18.9\%(정제29)에서 94.1\%(정제30)로 올라 두 장 모두 정답이 되었다. 갈색여치 15장 전체에 같은 처리를 하면 평균 정답 확률이 82.9\%에서 94.4\%로 오르고 15장이 모두 정답이 되었다. 두 장 모두 몸통이 밝게 찍힌 것이 원인이라는 짐작이 맞았다.

다리의 무늬와 가시가 흐림 (natureobserve). 다리가 흐려서 모델이 다리를 단서로 쓰지 못했다면, 다리를 지워도 정답 확률은 거의 변하지 않아야 한다. 그런데 natureobserve의 다리를 지우자 18.9\%(정제29)에서 8.2\%(정제31)로 더 떨어졌다. 흐린 다리에서도 모델은 갈색여치 단서를 얻고 있었으므로, 다리가 흐린 것은 오답의 주된 원인이 아니었다.

\begin{center}\footnotesize
\begin{tabular}{lll}
\toprule
짐작한 원인 & 검증 방법 & 결과 \\
\midrule
몸통이 밝게 찍힘 (두 장) & 몸통 바탕 검게 & add\_031 34.7\% → 67.9\%, natureobserve 18.9\% → 94.1\% → 원인 맞음 \\
아랫배 초록이 안 보임 (두 장) & 초록 아랫배 제거 (가설 4) & 영향 없음 → 원인 아님 \\
다리 무늬·가시가 흐림 (natureobserve) & 다리 제거 & 18.9\% → 8.2\%, 더 떨어짐 → 주된 원인 아님 \\
\bottomrule
\end{tabular}
\end{center}

\sidebyside{figures/p7_023.jpg}{정제23 test/add\_031 원본: 34.7\% (오답)}
           {figures/p7_028.jpg}{정제28 몸통 바탕 검게: 67.9\%}
\sidebyside{figures/p7_029.jpg}{정제29 test/natureobserve 원본: 18.9\% (오답)}
           {figures/p7_030.jpg}{정제30 몸통 바탕 검게: 94.1\%}
\sidebyside{figures/p7_029.jpg}{정제29 test/natureobserve 원본: 18.9\% (오답)}
           {figures/p7_031.jpg}{정제31 다리 제거: 8.2\% (오답)}

\subsection*{가설 6. 갈색여치 분류도 배경의 영향을 받는다}

가설 5까지는 곤충 자체의 요소만 보았다. 하지만 좀날개여치 분석의 가설 3에서, 풀잎 배경 갈색여치 4장의 배경을 검정으로 바꾸자 3장이 오답이 되었다. 갈색여치 분류도 배경의 영향을 받는다는 것이 이미 확인된 것이다. 그래서 배경을 하나의 요소로 추가해 갈색여치 15장 전체에서 다시 한번 검증하고, 배경이 몸통·가슴·다리보다 강한지 교차해 비교했다.

곤충은 그대로 두고 배경만 흰색, 검정, 풀잎으로 바꿨다. 풀잎 배경은 다른 갈색여치 사진(test/추가\_0038)에서 곤충이 없는 부분만 잘라 썼다. 배경만 바꿨을 때의 결과는 다음과 같다.

\begin{center}\footnotesize
\begin{tabular}{lrr}
\toprule
배경 & 평균 정답 확률 (정답이던 13장) & 15장 중 정답 \\
\midrule
원본 배경 & 82.9\% & 13/15 \\
흰 배경 & 70.5\% & 13/15 \\
검정 배경 & 12.2\% & 1/15 \\
풀잎 배경 & 85.1\% & 15/15 \\
\bottomrule
\end{tabular}
\end{center}

검정 배경에서는 15장 중 14장이 오답이 되었다. 가장 확실하게 정답이던 new\_images/p3도 98.6\%(정제38)에서 30.3\%(정제39)로 떨어졌다. 반대로 풀잎 배경에서는 15장이 모두 정답이었고, 원래 오답이던 두 장도 배경만 풀잎이나 흰색으로 바꾸면 정답이 되었다(정제32--37).

\par\noindent
\pic{figures/p7_023.jpg}{정제23 add\_031 원본: 34.7\%}\hfill
\pic{figures/p7_032.jpg}{정제32 흰 배경: 97.1\%}\hfill
\pic{figures/p7_033.jpg}{정제33 검정 배경: 15.9\%}\hfill
\pic{figures/p7_034.jpg}{정제34 풀잎 배경: 99.3\%}
\par\medskip
\par\noindent
\pic{figures/p7_029.jpg}{정제29 natureobserve 원본: 18.9\%}\hfill
\pic{figures/p7_035.jpg}{정제35 흰 배경: 81.8\%}\hfill
\pic{figures/p7_036.jpg}{정제36 검정 배경: 0.5\%}\hfill
\pic{figures/p7_037.jpg}{정제37 풀잎 배경: 96.7\%}
\par\medskip

\sidebyside{figures/p7_038.jpg}{정제38 new\_images/p3 원본: 98.6\%}
           {figures/p7_039.jpg}{정제39 검정 배경: 30.3\% (오답, 좀날개여치)}

다음으로 배경의 영향이 몸통·가슴·다리보다 강한지 알기 위해, 배경 4가지와 처리 4가지(그대로, 가슴 가림, 다리 제거, 몸통 바탕 검게)를 모두 조합해 15장에 적용했다.

\begin{center}\footnotesize
\begin{tabular}{lrrrr}
\toprule
배경 \textbackslash{} 처리 & 그대로 & 가슴 가림 & 다리 제거 & 몸통 바탕 검게 \\
\midrule
원본 배경 & 82.9\% (13/15) & 71.2\% (12/15) & 47.4\% (8/15) & 94.4\% (15/15) \\
흰 배경 & 70.5\% (13/15) & 58.8\% (13/15) & 26.1\% (3/15) & 88.9\% (15/15) \\
검정 배경 & 12.2\% (1/15) & 5.7\% (0/15) & 9.2\% (0/15) & 21.1\% (2/15) \\
풀잎 배경 & 85.1\% (15/15) & 75.8\% (14/15) & 44.4\% (7/15) & 97.7\% (15/15) \\
\bottomrule
\end{tabular}
\end{center}
각 칸은 원래 정답이던 13장의 평균 정답 확률이고, 괄호는 15장 중 정답 수다.

검정 배경은 모든 처리를 이겼다. 다른 배경에서는 15장을 모두 정답으로 만든 `몸통 바탕 검게'도 검정 배경에서는 2장만 정답이었다. natureobserve는 풀잎 배경에서 몸통 바탕을 검게 하면 99.6\%(정제43)지만, 검정 배경에서는 5.2\%(정제42)였다.

다리는 배경과 별개로 중요했다. 풀잎 배경에서 다리를 지우면 7/15로, 원본 배경에서 다리를 지운 경우(8/15)와 비슷했다. 흰 배경에서 다리를 지우면 3/15로 더 낮아졌다. test/add\_013은 흰 배경에서 다리를 지우자 90.2\%(정제20)에서 27.9\%(정제44)로 떨어졌다.

검은 몸통은 검정 배경만 아니면 어디서나 강했다. 몸통 바탕을 검게 하면 원본·흰색·풀잎 배경에서 모두 15장이 정답이었다.

가슴 무늬는 어느 배경에서나 영향이 작았다.

즉 영향의 크기는 검정 배경이 가장 컸고, 다리와 검은 몸통이 그다음, 가슴 무늬가 가장 작았다. 풀잎 배경은 정답 확률을 조금 올려 주었다.

오답 두 장의 교차 결과는 다음과 같다.

\begin{center}\footnotesize
\begin{tabular}{lrrrr}
\toprule
test/add\_031 & 그대로 & 가슴 가림 & 다리 제거 & 몸통 바탕 검게 \\
\midrule
원본 배경 & 34.7\% (오답) & 54.4\% & 13.8\% (오답) & 67.9\% \\
흰 배경 & 97.1\% & 94.3\% & 54.1\% & 98.4\% \\
검정 배경 & 15.9\% (오답) & 21.2\% (오답) & 5.7\% (오답) & 25.1\% (오답) \\
풀잎 배경 & 99.3\% & 98.8\% & 93.4\% & 99.6\% \\
\midrule
test/natureobserve & 그대로 & 가슴 가림 & 다리 제거 & 몸통 바탕 검게 \\
\midrule
원본 배경 & 18.9\% (오답) & 4.9\% (오답) & 8.2\% (오답) & 94.1\% \\
흰 배경 & 81.8\% & 47.7\% & 18.9\% (오답) & 95.9\% \\
검정 배경 & 0.5\% (오답) & 0.4\% (오답) & 2.7\% (오답) & 5.2\% (오답) \\
풀잎 배경 & 96.7\% & 76.6\% & 18.2\% (오답) & 99.6\% \\
\bottomrule
\end{tabular}
\end{center}

두 장의 차이는 다리에서 나타났다. add\_031은 풀잎 배경이면 다리를 지워도 93.4\%(정제40)로, 이 사진에서는 배경이 다리보다 강했다. 반면 natureobserve는 풀잎 배경이어도 다리를 지우면 18.2\%(정제41)로 떨어져, 다리가 배경보다 강했다. 경계 사례인 add\_008은 풀잎 배경에서 84.6\%, 추가\_0039는 몸통 바탕을 검게 했을 때 80.1\%로 확실한 정답이 되었다.

\par\noindent
\pic[0.31]{figures/p7_040.jpg}{정제40 add\_031 풀잎 + 다리 제거: 93.4\%}\hfill
\pic[0.31]{figures/p7_041.jpg}{정제41 natureobserve 풀잎 + 다리 제거: 18.2\%}\hfill
\pic[0.31]{figures/p7_042.jpg}{정제42 natureobserve 검정 + 몸통 바탕 검게: 5.2\%}
\par\medskip
\sidebyside{figures/p7_043.jpg}{정제43 natureobserve 풀잎 + 몸통 바탕 검게: 99.6\%}
           {figures/p7_044.jpg}{정제44 add\_013 흰 배경 + 다리 제거: 27.9\%}

학습 데이터의 배경을 종별로 세어 보았다(\texttt{train\_backgrounds.py}). 풀잎 배경 비율은 갈색여치가 39\%(59장 중 23장)로 5종 중 가장 높았고, 좀날개여치는 23\%였다. 어두운 배경은 갈색여치가 0장이었다. 풀잎 배경이 갈색여치 판단을 돕고 어두운 배경이 방해하는 것은 이 분포와 맞는다. 다만 회색 돌·콘크리트 배경 비율은 갈색여치 44\%, 좀날개여치 47\%로 비슷해서, 파이프·돌 배경의 어떤 점이 분류를 끌어내렸는지까지는 밝히지 못했다.

% ----- 여치
\subsection*{여치 분석 개요}

테스트와 new\_images의 여치(Gampsocleis sedakovii obscura) 15장 중 11장은 정답, 4장은 오답이었다. test/add\_021(정제45), test/add\_028(정제46), test/다운로드 (13)(정제47)은 긴날개여치로, new\_images/p5(정제48)는 중베짱이로 분류되었다.

\par\noindent
\pic{figures/p7_045.jpg}{정제45 add\_021: 36.4\%}\hfill
\pic{figures/p7_046.jpg}{정제46 add\_028: 6.7\%}\hfill
\pic{figures/p7_047.jpg}{정제47 다운로드 (13): 17.1\%}\hfill
\pic{figures/p7_048.jpg}{정제48 p5: 29.3\%}
\par\medskip

나는 여치를 다음 특징으로 구별한다고 생각했다.
\begin{enumerate}
\item 배의 얼룩: 배에 초록과 검정의 얼룩이 심하다.
\item 날개의 얼룩: 날개에 초록과 갈색의 얼룩이 있고, 날개 가장자리는 갈색이다.
\item 두 가지 발색: 여치에는 초록 발색과 갈색 발색이 나타난다. 따라서 분류기가 두 부류로 학습했을 수 있지만, 두 발색 모두 배와 날개의 얼룩이 심하므로 결국 이 얼룩으로 여치를 구별할 가능성이 크다.
\end{enumerate}

오답 이미지를 보면, add\_028을 빼고 세 장(add\_021, 다운로드 (13), p5)은 모두 배가 다리에 가려져 배의 얼룩이 보이지 않는다는 공통점이 있다. 나는 이것이 오답의 원인이라고 생각했다. add\_028은 배가 가려지지 않았는데도 틀려 원인을 알 수 없었다. 다만 긴날개여치로 분류된 점과 회색 아스팔트 위에서 찍힌 점을 보고, 긴날개여치 학습 사진에 회색 배경이 많기 때문일지도 모른다고 생각했다.

갈색여치 분석과 같은 순서로 검증했다. 먼저 모델이 배와 날개의 얼룩을 실제로 쓰는지 확인하고(가설 7), 그 얼룩을 발색과 상관없이 보는지 확인한 뒤(가설 8), 이를 바탕으로 오답 세 장(가설 9)과 add\_028(가설 10)을 검증했다. 배와 날개는 갈색여치의 가슴처럼, 사진마다 위치를 직접 정한 회색 네모로 가렸다. 수치는 모두 여치 정답 확률이다.

\textbf{정제 요청.} 배와 날개는 갈색여치 때처럼 네모로 가려 달라고 했고, 원인을 알 수 없는 add\_028은 배경을 바꾼 이미지도 요청했다.

에이전트에게 한 말 그대로:
\begin{quote}
"돼지여치의 배를 갈색여치 때처럼 네모 박스로 가려줘. 날개도 가리는 정제해줘. 그리고 add28은 왜 틀린건지 원인을 모르겠어. 아무래도 배경의 원인이 있는지도 몰라. 지금 대충 보니까 긴날개여치 배경에 회색 바닥이 많은 것 같거든? add28 배경을 다양하게 바꾸는 정제해줘. 풀잎 배경도 있으면 좋겠어."
\end{quote}

\subsection*{가설 7. 모델은 배와 날개의 얼룩으로 여치를 구별한다}

배와 날개를 각각, 그리고 함께 회색 네모로 가려 15장에 적용했다. 배 네모는 머리 바로 뒤부터 배 끝까지, 날개 네모는 날개가 시작하는 곳부터 날개 끝까지 덮도록 사진마다 직접 정했다(\texttt{change\_gampsocleis.py}의 VIEW, 격자는 \texttt{grid\_for\_boxes.py}). 가렸을 때 정답 확률이 크게 떨어지면 모델이 그 부위의 얼룩을 쓰고 있는 것이다.

\textbf{정제 요청.} 여치 15장 각각에 대해 날개만, 배만, 둘 다 회색 네모로 가린 이미지를 요청했다. 배 네모는 머리 바로 뒤부터 배 끝까지, 날개 네모는 날개가 시작하는 곳부터 날개 끝까지 덮어야 한다. 처음 만든 네모는 머리 쪽으로 치우쳐 배 끝과 날개 끝이 남아서 두 네모를 여치의 뒤쪽으로 옮기게 했고, 배 네모가 다리를 덮은 add\_033은 다시 정하게 했다.

에이전트에게 한 말 그대로:
\begin{quote}
"배랑 날개 가리는 거 보면 배를 가리는건 머리 뒤부터 엉덩이까지 가려야하는데 지금은 머리 시작부터 엉덩이 좀 남게 가리고 있어. 날개도 마찬가지로 날개 끝을 못가리며 대신 머리까지 가리고 있어. 두 가림막 다 좀 더 여치의 뒤쪽으로 이동해야해. 너가 가림 필터 적용한 이미지 보면 무슨 말인지 알거야"
\end{quote}

\begin{center}\footnotesize
\begin{tabular}{lrrr}
\toprule
가린 부위 & 정답 확률 평균 변화 (정답이던 11장) & 오답으로 바뀜 & 15장 중 정답 \\
\midrule
날개 & $-$43.6\%p & 6/11 & 5/15 \\
배 & $-$53.5\%p & 7/11 & 4/15 \\
날개 + 배 & $-$71.7\%p & 11/11 & 0/15 \\
\bottomrule
\end{tabular}
\end{center}

배와 날개 모두 여치의 큰 단서였고, 배가 날개보다 조금 더 강했다. 둘을 함께 가리면 정답이던 11장이 모두 오답이 되어, 여치 판단은 거의 전부 이 두 부위에 달려 있었다. 가설 7은 맞았다.

예를 들어 test/add\_029는 원본 88.6\%(정제49)에서 날개를 가리면 24.8\%(정제50), 배를 가리면 3.3\%(정제51), 둘 다 가리면 2.0\%(정제52)가 되었다. new\_images/p3도 90.1\%(정제53)에서 각각 60.5\%(정제54), 20.9\%(정제55), 14.6\%(정제56)로 떨어졌다.

\par\noindent
\pic{figures/p7_049.jpg}{정제49 add\_029 원본: 88.6\%}\hfill
\pic{figures/p7_050.jpg}{정제50 날개 가림: 24.8\%}\hfill
\pic{figures/p7_051.jpg}{정제51 배 가림: 3.3\%}\hfill
\pic{figures/p7_052.jpg}{정제52 날개 + 배 가림: 2.0\%}
\par\medskip
\par\noindent
\pic{figures/p7_053.jpg}{정제53 p3 원본: 90.1\%}\hfill
\pic{figures/p7_054.jpg}{정제54 날개 가림: 60.5\%}\hfill
\pic{figures/p7_055.jpg}{정제55 배 가림: 20.9\%}\hfill
\pic{figures/p7_056.jpg}{정제56 날개 + 배 가림: 14.6\%}
\par\medskip

이미지별 전체 결과는 다음과 같다. 발색은 내가 눈으로 구분했고, 마지막 열은 가설 8의 결과다.

\begin{center}\footnotesize
\begin{tabular}{llrrrrr}
\toprule
이미지 & 발색 & 원본 & 날개 가림 & 배 가림 & 날개 + 배 가림 & 발색 바꿈 \\
\midrule
test/add\_011 & 갈색 & 86.5\% & 36.4\% & 67.1\% & 20.3\% (오답) & 69.5\% \\
test/add\_021 & 초록 & 36.4\% (오답) & 12.7\% (오답) & 6.0\% (오답) & 6.4\% (오답) & 67.5\% \\
test/add\_024 & 갈색 & 99.0\% & 84.9\% & 59.1\% & 15.4\% (오답) & 89.2\% \\
test/add\_028 & 갈색 & 6.7\% (오답) & 3.0\% (오답) & 1.1\% (오답) & 0.5\% (오답) & 0.3\% (오답) \\
test/add\_029 & 초록 & 88.6\% & 24.8\% (오답) & 3.3\% (오답) & 2.0\% (오답) & 81.3\% \\
test/add\_030 & 초록 & 92.4\% & 70.0\% & 32.4\% (오답) & 7.6\% (오답) & 87.5\% \\
test/add\_033 & 초록 & 91.5\% & 20.6\% (오답) & 3.9\% (오답) & 6.3\% (오답) & 96.1\% \\
test/다운로드 (11) & 초록 & 95.3\% & 5.2\% (오답) & 47.3\% & 8.2\% (오답) & 95.8\% \\
test/다운로드 (12) & 초록 & 47.1\% & 18.8\% (오답) & 2.2\% (오답) & 3.3\% (오답) & 71.1\% \\
test/다운로드 (13) & 초록 & 17.1\% (오답) & 11.3\% (오답) & 4.1\% (오답) & 6.5\% (오답) & 51.0\% \\
new/p1 & 갈색 & 85.7\% & 7.1\% (오답) & 44.0\% & 6.2\% (오답) & 27.0\% (오답) \\
new/p2 & 초록 & 53.7\% & 33.6\% (오답) & 15.9\% (오답) & 5.9\% (오답) & 90.2\% \\
new/p3 & 초록 & 90.1\% & 60.5\% & 20.9\% (오답) & 14.6\% (오답) & 89.5\% \\
new/p4 & 갈색 & 85.2\% & 73.3\% & 30.5\% (오답) & 36.9\% (오답) & 48.0\% \\
new/p5 & 초록 & 29.3\% (오답) & 0.7\% (오답) & 1.4\% (오답) & 0.5\% (오답) & 86.5\% \\
\bottomrule
\end{tabular}
\end{center}

\subsection*{가설 8. 모델이 보는 얼룩은 발색과 상관없이 같다}

가설 7로 모델이 배와 날개의 얼룩을 본다는 것은 확인했다. 하지만 이것만으로는 모델이 그 얼룩을 어떤 형태로 기억하는지는 알 수 없다. 여치에는 초록형과 갈색형이 있으므로, 두 가지 가능성이 있다.
\begin{enumerate}
\item 얼룩 자체만 보기보단 얼룩을 발색과 묶어서 기억한다. `초록 바탕의 검은 띠'와 `갈색 바탕의 검은 띠'를 따로 배웠다면, 발색이 바뀌는 순간 같은 얼룩도 알아보지 못하거나 그 가중치가 달라질 것이다.
\item 얼룩의 무늬만 본다. 발색과 상관없이 얼룩의 형태를 본다면, 발색을 바꿔도 정답이 유지되어야 한다.
\end{enumerate}
이 차이는 오답 분석에서 중요하다. 오답 세 장이 모두 초록형이므로, 발색 자체가 오답에 관여하는지도 여기서 함께 확인할 수 있다. 학습 사진 39장을 눈으로 나눠 보니 초록형이 약 24장, 갈색형(적갈색·보라빛 포함)이 약 15장으로 두 발색이 모두 충분히 들어 있었고, 테스트 15장은 초록형 10장, 갈색형 5장이었다.

두 가능성을 가르기 위해 곤충 안의 색상만 바꿔, 초록형은 갈색형으로, 갈색형은 초록형으로 만들었다. 밝기와 무늬는 그대로이므로 얼룩의 형태는 남고 발색만 바뀐다.

\textbf{정제.} 여치 15장에서 곤충 안의 색상만 바꿔, 초록형은 갈색형으로, 갈색형은 초록형으로 만든 이미지로 검증했다. 밝기와 무늬는 그대로 둔다.

에이전트에게 한 말 그대로:
\begin{quote}
"돼지여치의 갈색형은 초록형으로, 초록형은 갈색형으로 바꿔줘. 이때 너무 발색을 완벽하게 바꿀 필요는 없어. 배와 날개의 무늬 경계는 다 살려야해. 색은 갈색형과 초록형에 가깝게만 바꿔도 돼."
\end{quote}

\begin{center}\footnotesize
\begin{tabular}{lrrr}
\toprule
처리 & 정답 확률 평균 변화 (정답이던 11장) & 오답으로 바뀜 & 15장 중 정답 \\
\midrule
발색 바꿈 & $-$6.4\%p & 1/11 & 13/15 \\
\bottomrule
\end{tabular}
\end{center}

발색을 바꿔도 정답이던 11장 중 10장이 정답을 유지했다. 따라서 모델은 얼룩을 발색과 묶어서 기억하기보다 얼룩의 무늬 자체를 보는 쪽에 가깝고, 가설 8은 대체로 맞았다.

다만 두 발색이 똑같이 강하지는 않았다. 원본에서 갈색형은 5장 중 4장, 초록형은 10장 중 7장이 정답이었다. 오답이던 초록형 세 장을 갈색형으로 바꾸자 모두 정답이 되었다. add\_021은 36.4\%(정제45)에서 67.5\%(정제60)로, p5는 29.3\%(정제48)에서 86.5\%(정제59)로, 다운로드 (13)은 17.1\%에서 51.0\%로 올랐다. 반대로 갈색형을 초록형으로 바꾸면 약해졌고, new\_images/p1은 85.7\%(정제57)에서 27.0\%(정제58)로 중베짱이가 되었다.

초록형 여치가 더 약한 이유는 긴날개여치와 중베짱이가 모두 초록색이기 때문으로 보인다. 갈색형은 색만으로도 두 종과 구별되지만, 초록형은 색이 같아서 배와 날개의 얼룩에 더 크게 기대야 한다.

\sidebyside{figures/p7_057.jpg}{정제57 p1 원본 (갈색형): 85.7\%}
           {figures/p7_058.jpg}{정제58 초록형으로: 27.0\% (오답, 중베짱이)}
\sidebyside{figures/p7_048.jpg}{정제48 p5 원본 (초록형): 29.3\% (오답)}
           {figures/p7_059.jpg}{정제59 갈색형으로: 86.5\%}
\sidebyside{figures/p7_045.jpg}{정제45 add\_021 원본 (초록형): 36.4\% (오답)}
           {figures/p7_060.jpg}{정제60 갈색형으로: 67.5\%}

\subsection*{가설 9. 오답 세 장은 배가 다리에 가려져서 틀렸다}

오답 세 장(add\_021, 다운로드 (13), p5)은 배가 다리에 가려져 있어서, 가려진 배를 드러낼 수는 없다. 대신 반대 방향으로 검증했다. 배가 보이지 않는 것이 원인이라면, 정답이던 사진의 배를 가렸을 때 오답 세 장과 같은 방식으로 틀려야 한다.

가설 7에서 배를 가리자 정답이던 11장 중 7장이 오답이 되었다. 그 7장은 모두 긴날개여치(3장)나 중베짱이(4장)로 분류되었는데, 이는 오답 세 장이 틀린 방식(긴날개여치 2장, 중베짱이 1장)과 같다. 예를 들어 test/add\_033은 배를 가리자 91.5\%(정제61)에서 3.9\%(정제62)로 떨어져 긴날개여치가 되었다.

또 배를 가렸을 때 초록형은 정답이던 7장 중 6장이 틀렸지만, 갈색형은 4장 중 1장만 틀렸다. 오답 세 장도 모두 초록형이다. 즉 초록형 여치는 배의 얼룩으로 초록색의 다른 종과 구별되는데, 배가 다리에 가려지면 이 단서를 잃고 긴날개여치나 중베짱이로 분류된 것이다. 가설 9는 맞았고, 가설 8에서 초록형을 갈색형으로 바꾸자 세 장 모두 정답이 된 것도 같은 이유로 설명된다.

\sidebyside{figures/p7_061.jpg}{정제61 add\_033 원본: 91.5\%}
           {figures/p7_062.jpg}{정제62 배 가림: 3.9\% (오답, 긴날개여치)}

\subsection*{가설 10. add\_028은 회색 배경 때문에 긴날개여치로 분류되었다}

add\_028은 배가 드러나 있는데도 여치 확률이 6.7\%밖에 되지 않았다. 배와 날개를 가리거나 발색을 바꿔도 모두 오답이었으므로, 곤충의 특징으로는 설명되지 않았다. 학습 데이터에서 긴날개여치 사진은 43장 중 34장(79\%)이 회색 돌·콘크리트 배경이었다. 그래서 회색 아스팔트 배경이 긴날개여치 판단을 만들었는지 확인하려고, 배경을 흰색, 검정, 단색 회색, 풀잎으로 바꿔 넣었다. 이번에는 여치 확률과 함께 긴날개여치와 1위 종의 확률도 적었다.

\begin{center}\footnotesize
\begin{tabular}{lrrl}
\toprule
add\_028 & 여치 & 긴날개여치 & 1위 \\
\midrule
원본 (회색 아스팔트, 정제46) & 6.7\% & 62.1\% & 긴날개여치 \\
흰 배경 (정제63) & 11.3\% & 0.4\% & 갈색여치 82.9\% \\
검정 배경 (정제64) & 10.8\% & 8.7\% & 좀날개여치 70.5\% \\
단색 회색 배경 (정제65) & 22.1\% & 6.2\% & 갈색여치 60.8\% \\
풀잎 배경 (정제66) & 0.3\% & 0.0\% & 갈색여치 97.9\% \\
곤충 크게 자름 (정제67) & 25.6\% & 40.2\% & 긴날개여치 \\
\bottomrule
\end{tabular}
\end{center}

배경을 바꾸자 긴날개여치 확률은 62.1\%에서 0.0--8.7\%로 사라졌다. 단색 회색 배경에서도 6.2\%였으므로, 긴날개여치 판단을 만든 것은 회색 그 자체가 아니라 아스팔트의 질감이다. 또 이 사진은 곤충이 화면의 4.0\%만 차지할 만큼 작게 찍혔다. 곤충을 크게 잘라 배경의 비중을 줄이자 여치 확률이 6.7\%에서 25.6\%로 오르고 긴날개여치는 40.2\%로 낮아졌다. 따라서 `긴날개여치'라는 오답은 배경 때문이라는 가설은 맞았다.

하지만 배경을 없애도 여치가 되지는 않았고, 대부분 갈색여치가 되었다. add\_028은 몸이 어두운 적갈색이다. 갈색여치 분석(가설 4)에서 모델은 어두운 몸통을 갈색여치의 큰 단서로 쓴다는 것을 확인했다. 그래서 배경의 영향이 사라지자 이 어두운 몸이 갈색여치로 읽힌 것으로 보인다. 여치 학습 사진 중 이처럼 몸이 어두운 적갈색인 개체는 몇 장뿐이었다.

\par\noindent
\pic[0.31]{figures/p7_046.jpg}{정제46 원본: 6.7\%}\hfill
\pic[0.31]{figures/p7_063.jpg}{정제63 흰 배경: 11.3\%}\hfill
\pic[0.31]{figures/p7_064.jpg}{정제64 검정 배경: 10.8\%}
\par\medskip
\par\noindent
\pic[0.31]{figures/p7_065.jpg}{정제65 단색 회색 배경: 22.1\%}\hfill
\pic[0.31]{figures/p7_066.jpg}{정제66 풀잎 배경: 0.3\%}\hfill
\pic[0.31]{figures/p7_067.jpg}{정제67 곤충 크게 자름: 25.6\%}
\par\medskip

% ----- 중베짱이
\subsection*{중베짱이 분석 개요}

테스트와 new\_images의 중베짱이(Tettigonia ussuriana) 15장 중 11장은 정답, 4장은 오답이었다. test/add\_017(정제68), new\_images/p1(정제69), new\_images/p5(정제71)는 긴날개여치로, new\_images/p4(정제70)는 좀날개여치로 분류되었다.

\par\noindent
\pic{figures/p7_068.jpg}{정제68 add\_017: 14.5\%}\hfill
\pic{figures/p7_069.jpg}{정제69 p1: 36.2\%}\hfill
\pic{figures/p7_070.jpg}{정제70 p4: 14.7\%}\hfill
\pic{figures/p7_071.jpg}{정제71 p5: 17.8\%}
\par\medskip

나는 중베짱이를 다음 특징으로 구별한다고 생각했다.
\begin{enumerate}
\item 밝기: 몸 색이 밝다.
\item 몸을 감싼 날개: 넓은 녹색 날개가 배를 포개어 몸 전체를 덮은 느낌이다.
\item 상단 갈색선: 전체적으로 무늬가 없지만, 날개 윗부분 가장자리만 갈색이고 이 갈색이 머리 위까지 이어진다.
\item 무늬 없는 배와 날개: 배와 날개 모두 다섯 종 중 가장 무늬가 없다.
\end{enumerate}
특히 긴날개여치와는 밝기, 날개, 그리고 상단 갈색선으로 구분한다고 보았다. 긴날개여치는 날개에만 갈색 무늬가 있고, 중베짱이는 머리까지 이어지기 때문이다.

오답을 보면 p1을 빼고 세 장은 모두 위에서 찍은 사진이다. 학습 데이터에 위에서 찍은 사진이 거의 없어 학습되지 않은 영역이라 틀렸다고 짐작했다. p1은 빛이 별로 없는 곳에서 찍혀 곤충이 어둡고, 그래서 가장 비슷하게 생기고 색이 더 어두운 긴날개여치로 오해받은 것으로 보았다. 또 모델이 머리까지 이어진 갈색선을 포착했다면 위에서 찍은 사진도 구분할 수 있었을 것이므로, 모델이 이 점을 포착하지 못했다고 생각했다.

앞선 종들과 같은 순서로 먼저 모델이 특징을 쓰는지 확인하고(가설 11), 그다음 p1의 밝기(가설 12)와 위에서 찍은 사진(가설 13)을 검증했다. 수치는 모두 중베짱이 정답 확률이다.

\textbf{정제 요청.} 중베짱이 15장에 대해, 머리-가슴-날개로 이어지는 상단 갈색선을 긴 네모로 가린 이미지와 그 안의 갈색만 주변 초록색으로 칠한 이미지를 함께 요청하고, 결과가 더 좋은 쪽을 쓰기로 했다. 날개가 배 아래까지 감싼 아랫부분을 가린 이미지도 요청했고, 밝기를 바꾸는 방법은 에이전트에게 맡겼다.

에이전트에게 한 말 그대로:
\begin{quote}
"상단의 머리-가슴-날개로 이어지는 갈색선은 기다란 박스를 여치 상단 끝부분에 위치시켜서 가려줘. 아니면 중베짱이의 머리 위 갈색 무늬를 주변 초록색으로 칠해줘. 둘 중에 더 결과 ㅗ좋은걸로 사용해. 밝기는 알아서 잘할거라고 생각하고, 날개가 몸을 감싼 형식은 배를 가려야할듯? 날개가 배 아래까지 감싸진 그 경계를 가려야해."
\end{quote}

\subsection*{가설 11. 모델은 상단 갈색선과 몸을 감싼 날개로 중베짱이를 구별한다}

두 특징을 따로 없애 15장에 적용했다(상자는 \texttt{boxes.py}).
\begin{itemize}
\item 상단 갈색선: 머리-가슴-날개 윗가장자리를 따라 긴 네모를 놓아 가리는 방법과, 그 네모 안의 갈색 픽셀만 주변 초록색으로 칠하는 방법을 모두 해 보았다. 네모로 가리면 갈색선 주변의 초록 부분까지 함께 사라지므로, 갈색선만 정확히 없애는 `칠하기'를 주 결과로 썼다. 위에서 찍은 사진은 등 가운데 줄무늬를 칠했다.
\item 몸을 감싼 날개: 날개가 배 아래까지 감싼 아랫부분을 네모로 가렸다. 옆모습 사진 11장에만 적용했다.
\end{itemize}

\begin{center}\footnotesize
\begin{tabular}{lrrr}
\toprule
없앤 특징 & 정답 확률 평균 변화 (정답이던 사진) & 오답으로 바뀜 & 정답 수 \\
\midrule
상단 갈색선 칠하기 & +0.7\%p & 0/11 & 15장 중 12장 (원본 11장) \\
(참고) 상단 네모 가림 & $-$7.3\%p & 1/11 & 15장 중 10장 \\
감싼 아랫부분 가림 & $-$26.0\%p & 3/10 & 옆모습 11장 중 7장 (원본 10장) \\
\bottomrule
\end{tabular}
\end{center}

상단 갈색선은 모델이 거의 쓰지 않았다. 갈색선을 초록으로 칠해도 정답 확률은 평균 0.7\%p밖에 변하지 않았다. test/add\_029는 칠하기 전후가 90.8\%(정제72)와 88.6\%(정제74)로 거의 같았고, 네모로 가렸을 때만 40.3\%(정제73)로 떨어졌다. 네모가 갈색선 주변의 날개까지 가렸기 때문이다. 즉 모델은 머리까지 이어진 갈색선을 중베짱이의 단서로 쓰지 않고 있었고, 내가 생각한 긴날개여치와의 구분법을 모델은 배우지 못했다.

오히려 위에서 찍은 new\_images/p4는 등의 주황·갈색 줄무늬를 초록으로 칠하자 14.7\%(정제70)에서 64.1\%(정제75)로 올라 정답이 되었다. 위에서 보면 갈색 줄무늬가 도드라져 모델이 갈색 계열인 좀날개여치로 본 것이다.

\par\noindent
\pic[0.31]{figures/p7_072.jpg}{정제72 add\_029 원본: 90.8\%}\hfill
\pic[0.31]{figures/p7_073.jpg}{정제73 상단 네모 가림: 40.3\%}\hfill
\pic[0.31]{figures/p7_074.jpg}{정제74 상단 갈색선 칠하기: 88.6\%}
\par\medskip
\sidebyside{figures/p7_070.jpg}{정제70 p4 원본 (위에서 찍음): 14.7\% (오답, 좀날개여치)}
           {figures/p7_075.jpg}{정제75 등 갈색 줄무늬 칠하기: 64.1\% (정답)}

몸을 감싼 날개는 큰 단서였다. 날개가 배를 감싼 아랫부분을 가리자 정답이던 10장 중 3장이 오답이 되었고, 세 장 모두 긴날개여치로 분류되었다. test/add\_012는 89.9\%(정제76)에서 34.7\%(정제77)로, new\_images/p2는 92.3\%(정제78)에서 36.8\%(정제79)로 떨어졌다. 날개가 몸을 감싼 모습이 보이지 않으면 중베짱이는 긴날개여치처럼 보인다는 뜻이다. 가설 11은 `몸을 감싼 날개'에 대해서만 맞았다.

\sidebyside{figures/p7_076.jpg}{정제76 add\_012 원본: 89.9\%}
           {figures/p7_077.jpg}{정제77 감싼 아랫부분 가림: 34.7\% (오답, 긴날개여치)}
\sidebyside{figures/p7_078.jpg}{정제78 p2 원본: 92.3\%}
           {figures/p7_079.jpg}{정제79 감싼 아랫부분 가림: 36.8\% (오답, 긴날개여치)}

\begin{center}\footnotesize
\begin{tabular}{llrrrrrr}
\toprule
이미지 & 촬영 & 원본 & 상단 칠하기 & 상단 네모 & 아랫부분 가림 & 이미지 어둡게 & 이미지 밝게 \\
\midrule
test/add\_012 & 옆 & 89.9\% & 81.3\% & 69.2\% & 34.7\% (오답) & 79.8\% & 91.2\% \\
test/add\_017 & 위 & 14.5\% (오답) & 18.5\% (오답) & 8.4\% (오답) & -- & 6.0\% (오답) & 17.9\% (오답) \\
test/add\_029 & 옆 & 90.8\% & 88.6\% & 40.3\% (오답) & 22.2\% (오답) & 87.7\% & 92.0\% \\
test/add\_030 & 옆 & 99.2\% & 99.1\% & 98.0\% & 75.7\% & 97.4\% & 99.5\% \\
test/다운로드 (12) & 옆 & 99.3\% & 99.3\% & 93.2\% & 87.7\% & 97.0\% & 99.8\% \\
test/다운로드 (13) & 옆 & 97.0\% & 97.5\% & 91.8\% & 76.0\% & 91.6\% & 96.2\% \\
test/다운로드 (15) & 옆 & 99.6\% & 98.7\% & 98.9\% & 87.6\% & 97.9\% & 99.5\% \\
test/다운로드 (19) & 옆 & 98.2\% & 98.3\% & 97.2\% & 96.1\% & 95.3\% & 98.8\% \\
test/다운로드 (23) & 옆 & 99.8\% & 99.8\% & 99.8\% & 95.6\% & 99.5\% & 99.8\% \\
test/다운로드 (8) & 옆 & 99.6\% & 99.7\% & 99.6\% & 93.5\% & 99.3\% & 99.7\% \\
new/p1 & 옆 & 36.2\% (오답) & 34.7\% (오답) & 17.3\% (오답) & 10.8\% (오답) & 14.1\% (오답) & 33.5\% (오답) \\
new/p2 & 옆 & 92.3\% & 90.3\% & 74.9\% & 36.8\% (오답) & 86.7\% & 94.5\% \\
new/p3 & 위 & 46.3\% & 67.5\% & 68.7\% & -- & 58.8\% & 38.0\% (오답) \\
new/p4 & 위 & 14.7\% (오답) & 64.1\% & 32.0\% (오답) & -- & 18.2\% (오답) & 28.8\% (오답) \\
new/p5 & 위 & 17.8\% (오답) & 18.4\% (오답) & 3.8\% (오답) & -- & 5.1\% (오답) & 43.8\% \\
\bottomrule
\end{tabular}
\end{center}

\subsection*{가설 12. p1은 곤충이 어둡게 찍혀서 긴날개여치로 분류되었다}

p1은 옆모습이고 날개가 몸을 감싼 모습도 보이는데 틀렸다. 먼저 사진 전체의 밝기를 1.3배로 밝게 해 보았지만 36.2\%에서 33.5\%(정제81)로 오히려 낮아졌다. p1을 다시 보면 흰 판 위에서 찍혀 사진 전체는 밝지만(밝기 143.6/255), 곤충만 따로 재면 78.1/255로 어둡다. 다른 중베짱이(옆모습 4장)는 106.1--187.2였다. 사진 전체를 밝게 하면 흰 배경도 함께 밝아져 곤충과의 대비가 그대로라서, 배경 제거 도구로 찾은 곤충 영역만 밝게 했다.

\begin{center}\footnotesize
\begin{tabular}{lrr}
\toprule
p1 처리 & 곤충 밝기 & 중베짱이 정답 확률 \\
\midrule
원본 (정제69) & 78.1 & 36.2\% (오답, 긴날개여치) \\
사진 전체 $\times$1.3 (정제81) & -- & 33.5\% (오답) \\
곤충만 $\times$1.3 & 100 & 67.1\% \\
곤충만 $\times$1.6 (정제80) & 123 & 84.0\% \\
곤충만 $\times$2.0 & 151 & 85.3\% \\
\bottomrule
\end{tabular}
\end{center}

곤충만 밝게 하자 다른 중베짱이 수준의 밝기에서 정답이 되었다. p1은 곤충이 어둡게 찍혀 긴날개여치로 분류된 것이 맞고, 가설 12는 맞았다. 배경을 흰색·검정·회색·풀잎으로 바꿔도 모두 오답(2.5--24.5\%)이었으므로 배경은 원인이 아니었다. 반면 정답이던 11장은 사진 전체를 어둡게($\times$0.7) 하거나 밝게($\times$1.3) 해도 평균 변화가 $-$1.9\%p, $-$0.3\%p로 거의 없었다. 밝기는 배경과 함께 바뀔 때가 아니라 곤충 자체의 밝기가 바뀔 때 영향을 주었다.

\par\noindent
\pic[0.31]{figures/p7_069.jpg}{정제69 p1 원본: 36.2\%}\hfill
\pic[0.31]{figures/p7_081.jpg}{정제81 사진 전체 밝게: 33.5\%}\hfill
\pic[0.31]{figures/p7_080.jpg}{정제80 곤충만 밝게: 84.0\%}
\par\medskip

\subsection*{가설 13. 위에서 찍은 사진은 학습 데이터에 거의 없어서 틀렸다}

촬영 방향은 이미지를 정제해서 바꿀 수 없으므로, 촬영 방향별 정답률과 학습 데이터의 구성을 비교했다.

\begin{center}\footnotesize
\begin{tabular}{lrl}
\toprule
중베짱이 테스트 & 사진 수 & 정답 \\
\midrule
옆에서 찍음 & 11장 & 10장 (오답은 p1) \\
위에서 찍음 & 4장 & 1장 (p3, 46.3\%로 간신히 정답) \\
\bottomrule
\end{tabular}
\end{center}

위에서 찍은 4장은 모두 틀리거나 경계에 걸렸고, 옆에서 찍은 사진은 p1만 틀렸다. 학습 사진을 눈으로 세어 보니 중베짱이 43장 중 위에서 찍은 사진은 약 4장뿐이었다. 반면 긴날개여치는 43장 중 약 7장이 위에서 찍은 사진이었다. 위에서 찍은 중베짱이 오답 3장 중 2장(add\_017, p5)이 긴날개여치로 분류된 것도, 위에서 본 모습을 긴날개여치 쪽으로 더 많이 배웠기 때문으로 보인다. 가설 13은 맞았다.

또 가설 11에서 본 것처럼, 위에서 찍으면 중베짱이의 큰 단서인 `몸을 감싼 날개'가 보이지 않고, 대신 등의 갈색 줄무늬가 도드라진다. p4처럼 이 줄무늬를 초록으로 칠하면 정답이 되는 경우도 있었다. 위에서 찍은 사진이 틀린 것은 학습 사진이 적은 데다, 모델이 쓰는 단서가 옆모습에서만 보이기 때문이다.

\begin{center}\pic[0.4]{figures/p7_082.jpg}{정제82 위에서 찍은 중베짱이 (p3): 46.3\%}\end{center}

% ----- 긴날개여치
\subsection*{긴날개여치 분석 개요}

테스트와 new\_images의 긴날개여치(Gampsocleis ussuriensis) 15장 중 8장은 정답, 7장은 오답이었고, 오답 7장은 모두 중베짱이로 분류되었다(정제83--89). 5종 중 오답이 가장 많다.

\par\noindent
\pic{figures/p7_083.jpg}{정제83 add\_012: 8.1\%}\hfill
\pic{figures/p7_084.jpg}{정제84 images (4): 43.3\%}\hfill
\pic{figures/p7_085.jpg}{정제85 img: 7.0\%}\hfill
\pic{figures/p7_086.jpg}{정제86 다운로드: 28.5\%}
\par\medskip
\par\noindent
\pic[0.31]{figures/p7_087.jpg}{정제87 p2: 33.1\%}\hfill
\pic[0.31]{figures/p7_088.jpg}{정제88 p4: 5.4\%}\hfill
\pic[0.31]{figures/p7_089.jpg}{정제89 p5: 22.9\%}
\par\medskip

나는 긴날개여치를 다음 특징으로 구별한다고 생각했다.
\begin{enumerate}
\item 가장 긴 날개: 다섯 종 중 날개가 가장 길어서, 배 끝 뒤로 날개가 얼마나 남는지가 중요하다. 날개는 배 위로 뻗어 있다.
\item 날개의 가로선 무늬: 날개에 가로선 얼룩이 조금 있다.
\end{enumerate}
또 여치 분석(가설 10)에서 회색 아스팔트 배경이 긴날개여치 판단으로 이어졌으므로, 긴날개여치도 배경의 영향을 받을 것이라고 보았다.

오답을 보면 대체로 밝은 빛 아래에서 찍힌 사진이 많아, 그 밝기 때문에 중베짱이로 분류된 것 같다고 짐작했다. 또 위에서 찍은 사진은 날개 길이와 가로선 무늬를 파악할 수 없어 오답을 냈을 것이라고 생각했다.

먼저 모델이 날개의 특징들을 쓰는지 확인하고(가설 14), 그 결과에서 이어진 배의 역할(가설 15), 그다음 배경과 밝기(가설 16)를 검증했다. 수치는 모두 긴날개여치 정답 확률이다.

\textbf{정제 요청.} 긴날개여치 15장에 대해 네 가지를 요청했다. 날개가 언제 끝나는지 모르도록 배 끝 뒤를 가린 이미지, 날개가 어디까지 내려오는지 모르도록 날개와 배의 경계를 가린 이미지, 날개의 진한 가로선 무늬만 주변 날개의 초록색으로 자연스럽게 칠한 이미지, 배경을 바꾼 이미지다. 밝기를 바꾸는 방법은 에이전트에게 맡겼다. 날개 무늬는 처음에 날개 전체를 흐리게 뭉개서 모자이크처럼 보였기 때문에, 진한 무늬만 골라 칠하도록 다시 만들게 했다.

에이전트에게 한 말 그대로:
\begin{quote}
"긴날개는 엉덩이 뒷부분부터 가리면 돼. 날개가 언제 끝나는지 모르도록. 그리고 밝기는 알아서 잘 할거고 날개의 배를 덮는 걸 검증하기 위해서는 배를 날개와 배의 경계를 가려서 날개가 어디까지 내려오는지 모르게해. 또 날개의 가로선 무늬를 체크하기 위해서 날개를 주변 초록색으로 칠해줘."
\end{quote}
\begin{quote}
"날개 무늬 지우는걸 좀 더 잘해야해. 지금은 그냥 모자이크 처리 되었는데 이게 아니고 이전에 갈색여치 무늬 없앤것처럼 날개의 진한 무늬들을 초록색으로 칠하는거지. 최대한 자연스럽게. 아 그리고 배경 바꾸는 정제해줘."
\end{quote}

\subsection*{가설 14. 모델은 날개 길이와 가로선 무늬로 긴날개여치를 구별한다}

세 가지 처리를 각각 따로 했다. 경계 가림은 옆모습 11장에만 적용했다.
\begin{itemize}
\item 날개 끝 가림: 배 끝부터 날개 끝까지를 네모로 가려, 날개가 언제 끝나는지 모르게 했다.
\item 날개-배 경계 가림: 날개의 아래 가장자리와 배가 만나는 띠를 가려, 날개가 배를 어디까지 덮는지 모르게 했다.
\item 날개 무늬 지우기: 날개 안에서 주변보다 확연히 어두운 가로선·얼룩 픽셀만 골라, 가까운 날개의 초록색으로 칠했다. 날개 본래의 색과 명암은 그대로 두었다.
\end{itemize}

\begin{center}\footnotesize
\begin{tabular}{lrrr}
\toprule
처리 & 정답 확률 평균 변화 (정답이던 사진) & 오답으로 바뀜 & 오답 → 정답 \\
\midrule
날개 끝 가림 & $-$1.9\%p & 1/8 & 0/7 \\
날개-배 경계 가림 & +9.4\%p & 0/4 & 3/7 \\
날개 무늬 지우기 & $-$18.5\%p & 2/8 & 0/7 \\
\bottomrule
\end{tabular}
\end{center}

날개 길이는 쓰이지 않았다. 배 뒤로 남는 날개를 가려도 정답 확률은 평균 1.9\%p만 떨어졌다. test/add\_003은 날개 끝을 가린 뒤에도 94.1\%(정제90)에서 95.5\%(정제91)로 그대로였다. 내가 가장 중요하다고 본 특징을 모델은 배우지 못했다.

\sidebyside{figures/p7_090.jpg}{정제90 add\_003 원본: 94.1\%}
           {figures/p7_091.jpg}{정제91 날개 끝 가림: 95.5\%}

날개와 배의 경계는 예상과 반대로 움직였다. 경계를 가리면 날개가 배를 어디까지 덮는지 알 수 없으므로, 긴날개여치의 특징이 사라져 중베짱이 확률이 오를 것이라고 예상했다. 그런데 중베짱이 확률은 오르지 않았고, 오히려 정답이던 4장은 긴날개여치 정답 확률이 평균 9.4\%p 올랐으며 오답 7장 중 3장은 정답이 되었다. test/다운로드는 28.5\%(정제86)에서 60.3\%(정제92)로, new\_images/p2는 33.1\%(정제87)에서 73.6\%(정제93)로 올랐다. 경계를 가려도 중베짱이 쪽으로 가지 않은 것은 예측 밖의 결과였다. 경계 띠에는 날개의 아래 가장자리와 함께 배의 윗부분이 들어 있으므로, 이는 배 자체에 단서가 있다는 의미일 수 있다. 그래서 다음 가설에서 배만 따로 가려 보았다.

\sidebyside{figures/p7_086.jpg}{정제86 다운로드 원본: 28.5\% (오답, 중베짱이)}
           {figures/p7_092.jpg}{정제92 날개-배 경계 가림: 60.3\% (정답)}
\sidebyside{figures/p7_087.jpg}{정제87 p2 원본: 33.1\% (오답, 중베짱이)}
           {figures/p7_093.jpg}{정제93 날개-배 경계 가림: 73.6\% (정답)}

가로선 무늬는 모델이 쓰는 단서였다. 날개의 진한 무늬를 초록으로 칠하자 정답이던 8장의 정답 확률이 평균 18.5\%p 떨어졌고, 2장은 중베짱이로 바뀌었다. test/add\_015는 89.8\%(정제94)에서 64.3\%(정제95)로, new\_images/p1은 65.3\%에서 16.0\%(정제96)로 떨어져 중베짱이가 되었다. 가로선 무늬가 없는 어두운 초록 날개는 무늬 없는 중베짱이 날개와 같아 보이기 때문이다. 다만 무늬를 지운다고 오답이 정답이 되지는 않았다. 가설 14는 가로선 무늬에 대해서는 맞았고, 날개 길이에 대해서는 틀렸다.

\par\noindent
\pic[0.31]{figures/p7_094.jpg}{정제94 add\_015 원본: 89.8\%}\hfill
\pic[0.31]{figures/p7_095.jpg}{정제95 날개 무늬 지우기: 64.3\%}\hfill
\pic[0.31]{figures/p7_096.jpg}{정제96 p1 날개 무늬 지우기: 65.3\% → 16.0\% (중베짱이)}
\par\medskip

\subsection*{가설 15. 긴날개여치를 구별하는 주요 단서는 배에 있다}

가설 14에서 날개와 배의 경계를 가리면 긴날개여치가 중베짱이 쪽으로 갈 것이라고 예상했지만, 오히려 긴날개여치 정답 확률이 올랐다. 경계 띠에는 날개의 아래 가장자리와 함께 배의 윗부분이 들어 있었으므로, 날개가 배를 덮는 정도가 아니라 배 자체에 긴날개여치의 주요 단서가 있을 수 있다고 생각했다. 이 가설이 맞다면, 배만 가렸을 때 긴날개여치 정답 확률은 떨어지고 중베짱이 쪽으로 가야 한다.

\textbf{정제 요청.} 긴날개여치 옆모습 11장에서 날개 아래로 드러난 배만 회색 네모로 가려 달라고 요청했다. 날개는 절대 가리지 않도록, 네모의 윗변은 날개의 아래 가장자리보다 아래에 두었다(상자는 \texttt{results/problem7/boxes.py}, 도구는 \texttt{change\_ussuriensis\_abdomen.py}).

에이전트에게 한 말 그대로:
\begin{quote}
"이를 위해서 정제해줘. 긴날개여치의 배만을 박스로 가려줘. 날개는 절대 가리면 안돼."
\end{quote}

\begin{center}\picturefile[0.9\linewidth]{figures/p7_abdomen_box_check.jpg}\end{center}
빨간 네모가 이번에 가린 배, 파란 네모가 가설 14에서 가린 날개-배 경계다.

\begin{center}\footnotesize
\begin{tabular}{lrrrr}
\toprule
처리 & 정답이던 4장 평균 변화 & 오답으로 바뀜 & 오답 7장 → 정답 & 11장 중 정답 \\
\midrule
배만 가림 & +9.7\%p & 0/4 & 2/7 & 6/11 \\
(참고) 날개-배 경계 가림 (가설 14) & +9.4\%p & 0/4 & 3/7 & 7/11 \\
\bottomrule
\end{tabular}
\end{center}

결과는 가설과 반대였다. 배만 가리자 11장 모두 긴날개여치 정답 확률이 올랐다. 정답이던 4장은 평균 9.7\%p 올랐고, 오답 7장 중 2장(test/images (4) 43.3\%→62.3\%, new\_images/p2 33.1\%→65.0\%)은 정답이 되었다. 정답이 되지 않은 오답들도 중베짱이 확률은 내려갔다(p4 74.7\%→59.9\%, img 72.7\%→61.9\%, 다운로드 68.6\%→57.3\%). test/add\_001은 45.3\%(정제106)에서 64.0\%(정제107)로 올랐다. 날개 아래로 드러난 배는 긴날개여치의 단서가 아니라 중베짱이 쪽 단서로 읽히고 있었으므로, 가설 15는 틀렸다.

배만 가렸을 때와 경계를 가렸을 때의 결과가 거의 같으므로(+9.7\%p와 +9.4\%p), 가설 14에서 경계 가림이 긴날개여치 확률을 올린 것도 날개의 가장자리보다는 그 아래의 배가 가려졌기 때문으로 보인다. 이는 중베짱이 분석(가설 11)에서 날개가 감싼 아랫부분을 가리면 긴날개여치로 바뀐 것과 같은 방향이다. 즉 모델은 두 종 모두에서 몸 아랫부분(배)이 보이는 모습을 중베짱이의 단서로 쓰고 있다. 

왜 이런 일이 생겼는지는 다음과 같이 추정한다. 분류기는 다섯 종이 확률을 나눠 갖기 때문에, 어떤 부분을 가리면 그 부분이 주던 증거가 빠지고 남은 종의 확률이 오른다. 배를 가렸을 때 중베짱이든 긴날개여치든 여치나 갈색여치가 아니라 긴날개여치 쪽으로 갔다는 것은, 초록색 여치에서 중베짱이의 증거(드러난 배)를 빼면 남는 첫 번째 후보가 긴날개여치라는 뜻이다. 반면 긴날개여치의 확률을 확실히 올린 것은 회색 배경, 어둡게 찍힌 곤충, 날개의 가로선 무늬처럼 몸 바깥이나 날개의 단서였다(가설 14, 16). 즉 모델은 긴날개여치를 `이런 몸이면 긴날개여치'로 배우기보다 `회색 배경 위의 어두운 초록 여치인데 중베짱이의 증거가 약하면 긴날개여치'로 배운 것으로 보인다. 학습 사진은 두 종 모두 43장이지만 긴날개여치는 79\%(34장)가 회색 돌·콘크리트 배경이어서, 몸보다 배경이라는 쉬운 단서로도 학습 사진을 맞힐 수 있었기 때문이라고 생각한다. 그래서 회색이 아닌 배경에서 밝게 찍힌 옆모습의 긴날개여치는 배가 보이는 순간 중베짱이로 읽혔고, 새 출처 사진 25장 중 모델이 긴날개여치라고 답한 것은 4장뿐이었다(\texttt{check.txt}). 학습 사진에서 긴날개여치의 배가 실제로 덜 보이는지는 확인하지 못했다.

\sidebyside{figures/p7_106.jpg}{정제106 add\_001 원본: 45.3\% (정답)}
           {figures/p7_107.jpg}{정제107 배만 가림: 64.0\%}
\sidebyside{figures/p7_087.jpg}{정제87 p2 원본: 33.1\% (오답, 중베짱이)}
           {figures/p7_108.jpg}{정제108 배만 가림: 65.0\% (정답)}

\subsection*{가설 16. 밝기와 배경이 긴날개여치 분류에 영향을 준다}

가설 14와 15의 날개와 배만으로는 오답 7장 중 최대 3장만 설명되었다. 그래서 곤충 바깥의 요소인 밝기와 배경을 검증했다.

밝기. 배경 제거 도구로 찾은 곤충 영역의 밝기를 재면, 오답 7장은 평균 134.0/255, 정답 8장은 120.5/255로 오답 쪽 곤충이 더 밝았다. 중베짱이 p1과 같은 방법으로 곤충만 어둡게($\times$0.75) 하거나 밝게($\times$1.3) 했다.

\begin{center}\footnotesize
\begin{tabular}{lrrrr}
\toprule
처리 & 정답이던 8장 평균 변화 & 오답으로 바뀜 & 오답 7장 → 정답 & 15장 중 정답 \\
\midrule
곤충만 어둡게 & +3.1\%p & 0/8 & 4/7 & 12/15 \\
곤충만 밝게 & $-$4.0\%p & 1/8 & 0/7 & 7/15 \\
\bottomrule
\end{tabular}
\end{center}

곤충을 어둡게 하자 오답 4장이 정답이 되었다. new\_images/p2는 33.1\%(정제87)에서 65.3\%(정제102)로, p5는 22.9\%(정제89)에서 44.1\%(정제103)로 올랐다. 중베짱이 p1(어두운 중베짱이 → 긴날개여치)과 반대 방향으로, 밝게 찍힌 긴날개여치는 중베짱이로 분류되었다. 모델은 두 종을 곤충의 밝기로도 나누고 있고, 밝은 빛이 오답의 원인이라는 짐작은 맞았다.

\sidebyside{figures/p7_102.jpg}{정제102 p2 곤충만 어둡게: 33.1\% → 65.3\%}
           {figures/p7_103.jpg}{정제103 p5 곤충만 어둡게: 22.9\% → 44.1\%}

배경. 배경을 제거한 뒤 흰색, 검정, 회색, 풀잎으로 바꿨다.

\begin{center}\footnotesize
\begin{tabular}{lrrrr}
\toprule
배경 & 정답이던 8장 평균 변화 & 오답으로 바뀜 & 오답 7장 → 정답 & 15장 중 정답 \\
\midrule
회색 & +0.8\%p & 0/8 & 5/7 & 13/15 \\
흰색 & $-$12.8\%p & 0/8 & 5/7 & 13/15 \\
검정 & $-$35.9\%p & 5/8 & 1/7 & 4/15 \\
풀잎 & $-$47.7\%p & 4/8 & 3/7 & 7/15 \\
\bottomrule
\end{tabular}
\end{center}

회색 배경에서는 정답이던 8장이 그대로이고 오답 5장이 정답이 되어, 15장 중 13장이 맞았다. test/img는 원래 식물 위에서 7.0\%(정제85)였지만 회색 배경에서 58.3\%(정제97)가 되었고, 흰 배경 30.2\%(정제98), 검정 배경 11.9\%(정제99), 풀잎 배경 57.7\%(정제100)였다. new\_images/p4도 자갈 위 5.4\%(정제88)에서 회색 배경 83.8\%(정제101)가 되었다. 반대로 검정·풀잎 배경에서는 정답 사진도 중베짱이나 갈색여치로 틀렸다.

긴날개여치 학습 사진은 43장 중 34장(79\%)이 회색 돌·콘크리트 배경이고, 풀잎 배경은 1장뿐이었다. 모델은 회색 배경을 긴날개여치의 단서로 배웠고, 회색이 아닌 배경(식물, 자갈, 어두운 바닥)에서는 긴날개여치를 알아보지 못했다. 가설 16은 밝기와 배경 모두에 대해 맞았다.

\par\noindent
\pic{figures/p7_097.jpg}{정제97 img 회색 배경: 58.3\%}\hfill
\pic{figures/p7_098.jpg}{정제98 img 흰 배경: 30.2\%}\hfill
\pic{figures/p7_099.jpg}{정제99 img 검정 배경: 11.9\%}\hfill
\pic{figures/p7_100.jpg}{정제100 img 풀잎 배경: 57.7\%}
\par\medskip
\sidebyside{figures/p7_088.jpg}{정제88 p4 원본 (자갈): 5.4\%}
           {figures/p7_101.jpg}{정제101 p4 회색 배경: 83.8\%}

\begin{center}\footnotesize
\begin{tabular}{llrrrrrrr}
\toprule
이미지 & 촬영 & 원본 & 날개 끝 & 경계 & 무늬 지움 & 곤충 어둡게 & 회색 배경 & 풀잎 배경 \\
\midrule
test/add\_001 & 옆 & 45.3\% & 44.9\% (오답) & 61.4\% & 36.1\% (오답) & 53.3\% & 88.3\% & 43.0\% \\
test/add\_003 & 옆 & 94.1\% & 95.5\% & 98.7\% & 57.3\% & 98.5\% & 91.5\% & 81.9\% \\
test/add\_012 & 옆 & 8.1\% (오답) & 9.4\% (오답) & 12.2\% (오답) & 7.6\% (오답) & 7.9\% (오답) & 69.7\% & 44.3\% \\
test/add\_015 & 옆 & 89.8\% & 81.5\% & 97.7\% & 64.3\% & 93.6\% & 81.3\% & 46.7\% \\
test/add\_026 & 위 & 90.5\% & 86.7\% & -- & 90.0\% & 73.3\% & 87.0\% & 11.7\% (오답) \\
test/images (4) & 옆 & 43.3\% (오답) & 40.3\% (오답) & 61.2\% & 32.9\% (오답) & 51.8\% & 61.0\% & 10.9\% (오답) \\
test/img & 옆 & 7.0\% (오답) & 8.2\% (오답) & 8.3\% (오답) & 1.6\% (오답) & 30.3\% (오답) & 58.3\% & 57.7\% \\
test/natureobserve & 위 & 87.5\% & 88.3\% & -- & 70.7\% & 95.1\% & 83.1\% & 8.3\% (오답) \\
test/다운로드 & 옆 & 28.5\% (오답) & 28.0\% (오답) & 60.3\% & 12.2\% (오답) & 49.2\% & 29.1\% (오답) & 21.2\% (오답) \\
test/추가\_0036 & 위 & 94.4\% & 93.2\% & -- & 93.0\% & 89.7\% & 85.8\% & 15.2\% (오답) \\
new/p1 & 옆 & 65.3\% & 64.4\% & 74.1\% & 16.0\% (오답) & 87.1\% & 72.4\% & 56.0\% \\
new/p2 & 옆 & 33.1\% (오답) & 11.1\% (오답) & 73.6\% & 17.9\% (오답) & 65.3\% & 60.3\% & 28.4\% (오답) \\
new/p3 & 위 & 96.9\% & 94.1\% & -- & 88.7\% & 98.1\% & 80.6\% & 19.3\% (오답) \\
new/p4 & 옆 & 5.4\% (오답) & 6.7\% (오답) & 7.5\% (오답) & 1.9\% (오답) & 9.1\% (오답) & 83.8\% & 45.2\% \\
new/p5 & 옆 & 22.9\% (오답) & 20.3\% (오답) & 28.2\% (오답) & 18.0\% (오답) & 44.1\% & 16.7\% (오답) & 1.7\% (오답) \\
\bottomrule
\end{tabular}
\end{center}

\subsection*{What I now think the model uses}

\textbf{좀날개여치.} 붉은 색감을 보정해도 오답은 고쳐지지 않았다. 풀잎 배경을 지우면 정답 확률이 12.9\%에서 52.5\%로 올라 겨우 정답이 되었지만, 같은 처리를 한 다른 풀잎 이미지들은 분류가 바뀌지 않았다. 반면 검정 배경은 모든 이미지를 좀날개여치 쪽으로 끌어당겼고, 이는 여치의 형태가 아닌 배경 밝기도 분류에 영향을 준다는 뜻이다.
따라서 이 오답은 하나의 잘못된 규칙이 아니라, 학습 데이터가 덮지 못한 개체에서 여러 약한 단서가 한쪽으로 누적된 결과로 해석된다. 배경 밝기처럼 종과 무관한 단서까지 학습된 것도 그 일부다. 좀날개여치 학습 사진이 43장뿐이라 개체 변이와 촬영 환경을 충분히 담지 못한 것이 근본 원인으로 추정된다.
개선하려면 좀날개여치 학습 사진을 늘리고, 어두운 배경의 갈색여치 사진도 넣어 배경 밝기가 한 종에만 몰리지 않게 해야 한다.
한계. 원인 분석은 오답 이미지 1장을 대상으로 했고, 가설 3을 확인한 풀잎 배경 이미지는 6장뿐이다. 배경을 제거한 뒤에도 일부 이미지에 작은 잎 조각이 남아 있어, 배경 외의 변화가 수치에 조금 섞였을 수 있다.

\textbf{갈색여치.} 내가 세운 특징 중 모델도 크게 쓰는 것은 검은 몸통과 다리(가시·마디)였다. 몸통을 갈색으로 바꾸면 정답 확률이 평균 33.0\%p, 다리를 지우면 35.4\%p 떨어졌다. 반면 가슴 무늬는 11.6\%p로 작은 단서였고, 초록 아랫배는 전혀 영향이 없었다.
오답 두 장에 대해서는, 몸통이 밝게 찍혔다는 짐작은 두 장 모두 맞았고, 아랫배 초록과 흐린 다리에 대한 짐작은 주된 원인이 아니었다. 또 몸통 바탕만 검게 해도, 배경만 풀잎이나 흰색으로 바꿔도 두 장 모두 정답이 되었다. 따라서 밝게 찍힌 몸통과 풀잎이 아닌 배경이 함께 오답을 만들었다.
교차 분석에서 가장 강한 요소는 검정 배경이었다. 검정 배경은 몸통을 검게 해도 되돌릴 수 없을 만큼 분류를 좀날개여치 쪽으로 끌어당겼고, 다리와 검은 몸통이 그다음, 가슴 무늬가 가장 약했다. 두 종의 분석을 합치면, 이 모델은 갈색여치와 좀날개여치를 곤충의 특징뿐 아니라 배경의 밝기와 종류로도 나누고 있다. 학습 데이터에서 풀잎 배경은 갈색여치에 많고 어두운 배경은 갈색여치에 없었기 때문에, 배경이 종을 구별하는 단서로 학습된 것으로 보인다. 개선하려면 풀잎이 아닌 배경과 어두운 배경의 갈색여치, 몸이 밝게 찍힌 갈색여치 사진을 학습 데이터에 더 넣어야 한다.
한계. 오답 이미지는 2장뿐이다. 초록 아랫배, 검은 몸통, 몸통 바탕 검게는 색과 밝기 기준으로 자동 처리해 경계가 완벽하지 않다. 가슴 네모는 사진마다 직접 정해 크기가 조금씩 다르고, 회색 네모 자체가 결과에 영향을 주었을 수 있다. 다리 제거는 더듬이도 함께 지웠고, 원본 배경에서는 일부 사진에 다리 그림자가 희미하게 남아 있다. 처음 시도한 `몸통 전체 밝기 낮추기'는 무늬까지 지워 결론을 잘못 이끌 뻔했으므로, 이미지 통제 방식 자체도 결과에 큰 영향을 준다.

\textbf{여치.} 내가 세운 특징대로 모델은 배와 날개의 얼룩으로 여치를 구별하고 있었다. 둘을 함께 가리면 정답 확률이 평균 71.7\%p 떨어져 정답이던 11장이 모두 틀렸고, 배가 날개보다 조금 더 강한 단서였다. 또 모델은 초록형과 갈색형을 모두 여치로 학습해, 발색을 바꿔도 대부분 정답을 유지했다.
오답 세 장은 배가 다리에 가려져서 틀렸다는 짐작이 맞았다. 정답이던 사진의 배를 가리면 같은 종(긴날개여치·중베짱이)으로 틀렸고, 이렇게 틀린 7장 중 6장이 초록형이었다. 초록형 여치는 색이 같은 긴날개여치·중베짱이와 배의 얼룩으로 구별되므로, 배가 가려지면 틀리기 쉽다.
add\_028의 `긴날개여치'라는 오답은 회색 아스팔트 배경 때문이었다. 곤충이 작게 찍혀 배경의 비중이 컸고, 학습 데이터의 긴날개여치는 79\%가 회색 배경이다. 다만 배경을 없애도 어두운 적갈색 몸이 갈색여치로 읽혀 여치가 되지는 않았다.
개선하려면 배가 가려진 자세의 초록형 여치, 작게 찍힌 여치, 몸이 어두운 적갈색 여치 사진을 학습 데이터에 더 넣어야 한다. 회색 배경이 긴날개여치에 몰리지 않도록 다른 종의 회색 배경 사진도 늘려야 한다.
한계. 오답 이미지는 4장뿐이다. 발색(초록형·갈색형)과 학습 사진의 발색 수는 내가 눈으로 구분한 것이다. 배·날개 네모는 사진마다 직접 정해 크기가 다르고, 네모가 다리 일부도 함께 가린 경우가 있다. 발색 바꾸기는 색상 기준으로 자동 처리해 일부 배경색도 함께 바뀔 수 있다.

\textbf{중베짱이.} 모델은 중베짱이를 날개가 몸을 감싼 모습과 곤충의 밝기로 구별하고 있었다. 감싼 아랫부분을 가리면 긴날개여치로 틀렸고, 곤충이 어두운 p1은 곤충만 밝게 하자 정답이 되었다. 반면 내가 긴날개여치와의 구분법으로 생각한 `머리까지 이어진 상단 갈색선'은 모델이 거의 쓰지 않았다.
오답 4장 중 3장은 위에서 찍은 사진이었다. 학습 데이터에 위에서 찍은 중베짱이가 약 4장뿐이고, 모델이 쓰는 `몸을 감싼 날개'는 위에서는 보이지 않는다. 남은 1장(p1)은 곤충이 어둡게 찍힌 것이 원인이었다.
개선하려면 위에서 찍은 중베짱이와 어두운 곳에서 찍힌 중베짱이 사진을 학습 데이터에 더 넣어야 한다.
한계. 오답은 4장이고, 위에서 찍은 사진의 수는 내가 눈으로 센 것이다. 상단 갈색선 칠하기는 네모 안의 갈색 계열 픽셀만 자동으로 골라 칠해 엷은 갈색이 일부 남았을 수 있다. `곤충만 밝게'는 p1에만 적용했다.

\textbf{긴날개여치.} 내가 가장 중요하다고 본 날개 길이는 모델이 쓰지 않았고, 날개의 가로선 무늬는 모델이 쓰는 단서였다(가설 14, 지우면 $-$18.5\%p). 날개와 배의 경계를 가리면 중베짱이 쪽으로 갈 것이라는 예상과 달리 긴날개여치 확률이 올라, 배에 긴날개여치의 주요 단서가 있다는 가설 15를 세웠다. 그러나 배만 가려도 11장 모두 긴날개여치 확률이 올라 이 가설은 틀렸고, 배는 오히려 중베짱이의 증거로 쓰이고 있었다. 정리하면 모델은 긴날개여치와 중베짱이를 다음으로 나누고 있었다.
\begin{itemize}
\item 날개의 가로선 무늬: 있으면 긴날개여치 쪽이다(가설 14).
\item 날개 아래로 드러난 배: 중베짱이의 증거다. 배를 가리면 어느 종이든 긴날개여치 쪽으로 기운다(가설 11, 15).
\item 곤충의 밝기: 밝으면 중베짱이, 어두우면 긴날개여치다(가설 16).
\item 배경: 회색 배경이면 긴날개여치 쪽이다(가설 16).
\end{itemize}
긴날개여치의 오답 7장은 모두 옆모습이었고 모두 중베짱이로 분류되었다. 곤충이 밝게 찍혔거나 회색이 아닌 배경에서 찍혔으며, 옆모습이라 중베짱이의 증거인 배가 그대로 보였다. 곤충을 어둡게 하거나(4장), 배경을 회색으로 바꾸거나(5장), 날개-배 경계를 가리거나(3장), 배만 가리면(2장) 정답이 되었다. 밝은 빛이 원인이라는 짐작은 맞았다.
중베짱이 분석과 합치면, 두 종은 같은 단서의 양쪽 끝에 놓여 있다. 모델은 두 종을 형태의 차이(날개 길이, 머리까지 이어진 갈색선)보다 밝기·배경·배가 보이는 모습으로 구분하고 있어, 이 중 하나만 어긋나도 서로 뒤바뀐다. 특히 긴날개여치는 몸의 특징보다 회색 배경, 어두운 몸, 날개 무늬처럼 몸 바깥이나 날개의 단서로 학습되어서, 이 단서가 약한 사진에서는 배가 보이는 순간 중베짱이로 읽힌다. 개선하려면 밝게 찍힌 긴날개여치, 회색이 아닌 배경의 긴날개여치, 어둡게 찍힌 중베짱이 사진을 학습 데이터에 더 넣어야 한다.
한계. 오답 7장 중 img, p4, add\_012는 배경을 바꾸는 것 외에는 어떤 처리로도 정답이 되지 않았다. 네모와 날개 영역은 사진마다 직접 정했고, 위에서 찍은 사진에는 경계 가림과 배 가림을 적용하지 않았다. 배 네모는 날개를 가리지 않도록 윗변을 날개의 아래 가장자리보다 낮게 두었기 때문에 날개 바로 아래의 배는 일부 남았고, 배를 가로지르는 다리가 함께 가려진 사진도 있다. 날개 무늬 지우기는 주변보다 20\% 이상 어두운 픽셀을 무늬로 보고 자동으로 칠했기 때문에, 날개 위에 걸친 다리의 어두운 부분도 일부 함께 칠해졌다. 촬영 방향과 학습 사진의 촬영 방향 수는 내가 눈으로 구분한 것이다.

\subsection*{Problem 7 outputs}

\texttt{results/problem7/measure.py}: 원본과 정제한 이미지 전부의 정답 확률. \texttt{x}는 오답이고, 그 뒤가 모델이 고른 종과 확률이다.
\outputfile{outputs/problem7_measure.txt}

\texttt{results/problem7/train\_backgrounds.py}: 학습 사진(data/clean)의 배경.
\outputfile{outputs/problem7_train_backgrounds.txt}

정제 도구가 출력한 상자 좌표와 밝기.
\outputfile{outputs/problem7_change_atlanticus.txt}
\outputfile{outputs/problem7_change_paratlanticus.txt}
\outputfile{outputs/problem7_change_gampsocleis.txt}
\outputfile{outputs/problem7_change_tettigonia.txt}
\outputfile{outputs/problem7_change_ussuriensis.txt}
\outputfile{outputs/problem7_change_ussuriensis_abdomen.txt}

% ---------------------------------------------------------------------------
\section*{Appendix. Output of check.py}
\outputfile{outputs/check.txt}

\end{document}
